\documentclass[11pt]{article}
\usepackage[margin=1in]{geometry}
\usepackage[T1]{fontenc}
\usepackage{lmodern,amsmath,amssymb,amsthm,booktabs,array,graphicx}
\usepackage[colorlinks=true,linkcolor=blue,citecolor=blue,urlcolor=blue]{hyperref}
\hypersetup{pdftitle={Depth dependence, internal robustness, and Bayesian neural-network learning},pdfauthor={ChatGPT-6 Astra; Claude Fable 5 and Opus 5}}
\usepackage{microtype}
\newtheorem{theorem}{Theorem}[section]
\newtheorem{proposition}[theorem]{Proposition}
\newtheorem{lemma}[theorem]{Lemma}
\newtheorem{corollary}[theorem]{Corollary}
\newtheorem{definition}[theorem]{Definition}
\newcommand{\E}{\mathbb E}
\newcommand{\Pp}{\mathbb P}
\newcommand{\R}{\mathbb R}
\newcommand{\X}{\{-1,1\}^d}
\newcommand{\op}{\mathrm{op}}
\newcommand{\Lip}{\operatorname{Lip}}
\newcommand{\sgn}{\operatorname{sign}}
\newcommand{\poly}{\operatorname{poly}}
\title{Depth dependence, internal robustness, and Bayesian neural-network learning}
\author{Research report prepared by ChatGPT-6 Astra and Claude (Fable 5 and Opus 5)\thanks{Authors are listed in alphabetical order.}\\for Kaarel H\"anni}
\date{17 October 2026}
\begin{document}
\maketitle
\begin{abstract}
We investigate whether the exponential depth dependence in the existing
width-independent Bayesian learning bounds is necessary, and which natural
robustness restrictions remove it. The known norm-capacity obstruction does
not establish an exponential sample lower bound for a small target under the
specified Gaussian prior; constants and single-coordinate targets in fact
have depth-independent sample bounds in the strong-prior limit.
We strengthen the capacity obstruction: even global
polynomial input-Lipschitz bounds allow all Boolean truth tables at polynomial
weighted norm and linear depth. Internal robustness is different. An
RMS-normalized all-layer margin $r$, together with energy $K$, permits
compression to width $O(Kr^{-2}(d+\log L))$. A second argument tracks entropy
only in unprocessed original weights while compressing a posterior. It yields
one polynomial-depth sample bound, under a single norm-calibrated robustness
restriction, simultaneously for every $\gamma\ge1$ and every width above the
width of a layered Boolean circuit. Neither a hard energy cutoff nor a
width-first limit is required. An alternative calibration uses the guaranteed
noise radius of a nonexpansive network with the same depth and readout.
Established all-layer-margin
and noise-stability literature supports the naturalness of the condition,
but not its universal validity for ordinary gradient-trained networks.
\end{abstract}

\paragraph{Scope and relation to later work.}
This report develops the mathematical arguments from the investigation of
17 October 2026. Its proofs are conventional and do not claim Lean
certification. We retain the all-input clean margin and posterior-confidence
conventions defined below. Fable's research note~\cite{fable}
supplies the robust-subsampling argument; this report derives its Bernstein
refinement, the Gaussian posterior bounds, and the capacity counterexamples.
The explicit bounds are deliberately loose and not proposed as practical
numerical estimates.
No claim of priority over the broader robustness/compression literature is made.

The source note's necessity discussion concerns uniform norm compression
and counting. Those properties alone do not imply a sample lower bound
for the Gaussian posterior, nor identify analog and robust target complexity
at the same architecture. Explicit robust circuit witnesses and posterior
arguments are therefore required below. As of 29 October 2026, a separate
result~\cite{conditional2026} rules out a universal sample bound polynomial
in target resources and learner depth, uniformly over all admissible learner
widths, assuming nonuniform $\mathsf P/\mathrm{poly}\ne\mathsf{TC}^1$.
That conditional obstruction supplements this report's historical open
question; the unconditional classification remains unresolved.

\section{Model and an all-layer margin}

Let $f:\X\to\{-1,1\}$ be a target and let $D$ be any probability
distribution on the Boolean cube. Training examples are independent
pairs $(X_i,f(X_i))$ with $X_i\sim D$. Let $L\ge2$, $d\ge1$,
and let every hidden layer have width $N$.
The input is $h^0=(1,x)$ for $x\in\X$; there are no subsequent biases.
Set
\[
 h^s=\phi(W^s h^{s-1})\quad(1\le s<L),\qquad
 z_w(x)=v h^{L-1},
\]
where $|\phi|\le1$ and $\phi$ is $1$-Lipschitz. Circuit constructions and
finite-variance target-mass estimates use $\phi=\tanh$. The weighted energy is
\[
 \kappa_N(w)=\frac dN\|W^1\|_F^2+
 \sum_{s=2}^{L-1}\|W^s\|_F^2+N\|v\|_2^2.
\]
For auxiliary classifiers produced by rounding, set $\sgn(0)=+1$.
Fix a prescribed scale schedule $t=t(d,L,N)>0$, written $t(N)$ when
$d,L$ are fixed; constant schedules are allowed. The raw entry variances
are $t^2/d,t^2/N,t^2/N^2$, respectively. With $\sigma^2=t^2/N$,
the Gaussian density in weighted coordinates is proportional
to $\exp(-\kappa_N/(2\sigma^2))$. A strong-prior limit sends
$\sigma\downarrow0$ at fixed $d,L,N$.
At exponent $\gamma$, set $t=N^{(1-\gamma)/2}$; the independent
entry variances are $N^{1-\gamma}/d$, $N^{-\gamma}$, and
$N^{-1-\gamma}$ for the first, internal, and readout layers respectively.
Without an extra support restriction, the posterior conditions the raw
Gaussian on $|z_w(x)|>1/2$ for every cube input and on
$f(X_i)z_w(X_i)>0$ for all observed inputs. A robustness-conditioned
posterior additionally imposes the support condition specified below.
All logarithms are natural unless a base is displayed.

Use additive \emph{postactivation} perturbations:
\[
 h^s_\eta=\phi(W^s h^{s-1}_\eta)+\eta^s,
 \qquad z_w(x;\eta)=v h^{L-1}_\eta+\eta^L.
\]
Their total normalized size is
\begin{equation}\label{eq:noisenorm}
 \|\eta\|_*^2=\sum_{s<L}\frac{\|\eta^s\|_2^2}{N}+|\eta^L|^2.
\end{equation}
The factor $N^{-1}$ measures hidden noise by its RMS amplitude and makes the
condition compatible with neuron cloning.

\begin{definition}
A network has all-layer sign margin $r$ for $f$ if
\[
 f(x)z_w(x;\eta)>0\quad\text{for every }x\in\X
 \text{ and every }\|\eta\|_*\le r.
\]
Let $\mathcal R_{r,N,f}$ additionally require the original clean margin
$f(x)z_w(x)>1/2$ for every $x$. The restricted prior is the raw Gaussian
conditioned on $\bigcup_f\mathcal R_{r,N,f}$, followed by conditioning on
the observed noiseless labels.
\end{definition}

The condition refers to a network's own clean signs; it does not reveal
unknown target labels to the prior. Its simultaneous finite-radius character
is important. A derivative evaluated only on the clean computation, random
noise tolerance, or robustness at one layer at a time is weaker.

This is an absolute-RMS version of the all-layer-margin idea of Wei and
Ma~\cite{weima2020}. Their normalization uses hidden-state norms, rather than
the fixed $\sqrt N$ used here. The two are not automatically equivalent when
hidden activity can collapse. Fable's note~\cite{fable} instead gives each layer
its own RMS budget $\tau$ and retains output margin $1/3$. The total norm and
maximum per-layer norm differ by at most $\sqrt L$; either convention can
give polynomial depth. The sharper proof below allocates noise by layer
energy and avoids even that extra factor in compressed width.

For a posterior $\Pi$, write
\[
 R_D(\Pi)=\E_{x\sim D}\Pp_{g\sim\Pi}\{g(x)\ne f(x)\}.
\]
The fraction of test inputs with correct-label probability at most $2/3$ is
at most $3R_D(\Pi)$. All learning bounds below therefore imply the original
confidence guarantee by replacing the desired error parameter by one third
of its value. They are not new majority-vote lower bounds.

\section{Findings and unresolved questions}

There are three distinct questions.
\begin{enumerate}
\item \emph{Can the old uniform norm-counting argument have polynomial depth?}
No. At depth $L=d+O(1)$ and width $2^{d+1}$, polynomial weighted norm already
contains every Boolean truth table. This remains true with a global input
Lipschitz constant polynomial in $d$.
\item \emph{Does the exact unrestricted Gaussian Bayesian learner therefore
need exponentially many samples for a polynomial-complexity target?}
This does not follow from norm counting, and was unresolved in the
17 October investigation. The missing
step is a target-specific prior/posterior mass or minimum-energy comparison.
\item \emph{Can a natural internal robustness condition give polynomial depth?}
Yes, under precise support and target hypotheses. An all-layer margin is an
established notion rather than a property invented solely for this proof.
Saturated Boolean circuits satisfy an inverse-polynomial version without
depth-dependent error accumulation.
\end{enumerate}

The resulting learning statements are summarized below. In the unified theorem,
``polynomial'' means one polynomial whose degree and coefficients are independent
of both depth and $\gamma\ge1$.

\begin{center}
\begin{tabular}{>{\raggedright\arraybackslash}p{.24\linewidth}>{\raggedright\arraybackslash}p{.37\linewidth}>{\raggedright\arraybackslash}p{.28\linewidth}}
\toprule
Regime & Additional support condition & Result established here\\
\midrule
Every $\gamma\ge1$, including mean field & Fixed all-layer margin, or either calibrated support rule in Sections~\ref{sec:selfcal} and~\ref{sec:readout-reference} & One polynomial bound; every width above circuit width\\[3pt]
Strong-prior limit & The same margin conditions & Polynomial depth; in fact every Gaussian scale $\sigma^2\le1/N$\\[3pt]
Any variance, including $\gamma=1$ & All-layer margin and a width-independent hard energy budget & Polynomial depth by a finite-class bound\\[3pt]
Unrestricted Gaussian prior & Original global clean margin only & General necessity unresolved; constants and dictators have a depth-independent strong-limit bound\\
\bottomrule
\end{tabular}
\end{center}

The circuit targets in this table are ordinary bounded-fan-in Boolean
circuits, expressed in layers. They are not silently identified with
arbitrary analog tanh networks at exactly the same depth. Polynomial-overhead
digital simulation gives an expressivity statement, but can increase both
depth and minimum required width; see Appendix~\ref{sec:digitize}.

\section{Compression with polynomial depth dependence}

\begin{theorem}\label{thm:compression}
Suppose $\kappa_N(w)\le K$, $K\ge1$, and $w$ has all-layer sign margin
$0<r\le1$ for $f$. Put
\[
 m=\left\lceil256Kr^{-2}[d+\log(400L)]\right\rceil.
\]
Then $f$ has a width-$m$, depth-$L$ implementation with clean margin $>3/4$,
energy at most $4K/r^2$, and every weight bounded in absolute value by $m$.
The compressed network need not itself retain the original robustness radius.
\end{theorem}

\begin{proof}
Assume first that $N\ge m$. For $2\le s<L$ put $k_s=\|W^s\|_F^2$,
and put $k_L=N\|v\|_2^2$, so $\sum_{s=2}^L k_s\le K$.
Zero-energy layers need no approximation. In each internal row, and in the
readout row, delete entries exceeding
\[
 \frac{8\sqrt K\,\|\text{row}\|_2}{r\sqrt N}
\]
in absolute value. Leave the first layer unchanged. Denote the pruned
weights by bars. Against any vector with coordinates in $[-1,1]$, the
internal pruning error has norm at most
$r\sqrt{Nk_s/K}/8$; the readout error is at most $r\sqrt{k_L/K}/8$.
Indeed, the removed row $\ell^1$ mass is at most its squared $\ell^2$ norm
divided by the deletion threshold.

Independently sample $m$ neurons with replacement from each hidden layer.
Keep sampled first-layer rows, and multiply all subsequent sampled source
weights by $N/m$. Define a virtual full-width forward pass that uses the
sampled sources at each layer. Its sampled coordinates are exactly the
small network's coordinates. Its layer-$s-1$ activation depends only on
samples through layer $s-2$, and is independent of the fresh source sample
used for layer $s$. This conditional independence is the central point in
the exact-simulation argument of Fable~\cite{fable}.

Fix an input and condition on those earlier samples. A fresh contribution
to an internal preactivation vector is
$Y=N\bar W^s_{\cdot J}\widetilde h^{s-1}_J$, with $J$ uniform on $[N]$.
Pruning and bounded activations give
\[
 \E\|Y-\E Y\|_2^2\le Nk_s,
 \qquad \|Y-\E Y\|_2\le\frac{16\sqrt{NKk_s}}r.
\]
For centered independent Hilbert-space variables of norm at most $B$ and
total second moment at most $V$, Bernstein's inequality gives
\[
 \Pp\!\left\{\left\|\sum_i Y_i\right\|\ge t\right\}
 \le2\exp\!\left(-\frac{t^2}{2(V+Bt/3)}\right).
\]
The dimension-free form used here is recorded, for example, in
Tarr\`es--Yao~\cite{tarres}, Appendix A, Corollary A.2, equation (A-1).
Applying it at sample-mean error
$r\sqrt{Nk_s/K}/8$ gives failure probability at most
\[
 2\exp\left(-\frac{3mr^2}{640K}\right)
 \le2\exp\left(-\frac{mr^2}{256K}\right).
\]
The readout calculation is identical with the factors $\sqrt N$ removed.
A union bound over all $2^d$ inputs and at most $L$ sampled layers costs
at most $1/200$ for the stated $m$.

On success, pruning plus sampling contributes normalized error at most
$r\sqrt{k_s/K}/4$ at layer $s$. Lipschitzness turns each preactivation
error into an additive postactivation error without increase. Thus the
virtual pass is an exact noisy pass of the original network with
\[
 \|\eta\|_*^2\le\sum_{s=2}^L\frac{r^2k_s}{16K}\le\frac{r^2}{16}.
\]
The expected weighted norm of the sampled network is at most $K$:
sampling is unbiased for every squared-norm term, with the prescribed
width rescalings. Markov's inequality gives norm at most $4K$ with
probability at least $3/4$. The good-noise and good-norm events consequently
intersect with probability at least $3/4-1/200$.

Every realized signed output in that intersection is $>3r/4$. Otherwise,
an additional scalar output perturbation of size at most $3r/4$ would
bring it to zero, with total noise norm at most $r$, contradicting the
all-layer sign margin. Scale the readout by $1/r$ to obtain margin $>3/4$
and energy at most $4K/r^2$.

If $N<m$, clone the original network to width $m$. Balanced multiplicities
increase weighted norm by at most two and preserve the output. Its clean
margin is $>r$, because output-only noise of size $r$ is allowed. Scaling
by $1/r$ gives the desired conclusion. Finally, $m\ge4K/r^2$ makes the
energy bound imply that each individual weight has magnitude at most $m$.
\end{proof}

Define $P=m(d+2)+(L-2)m^2$ and the conservative entropy bound
\begin{equation}\label{eq:entropy}
 H(d,L,K,r)=P\log\bigl[3+8m(d+L+1)(2m)^{2L}\bigr],
\end{equation}
with $m$ from Theorem~\ref{thm:compression}. A parameter grid of sufficiently
small mesh preserves the $>3/4$ margin: a safe output perturbation bound is
$\delta(d+L+1)(2m)^{2L}$ when each parameter is changed by at most $\delta$.
The number of functions under the theorem's hypotheses is therefore at most
$\exp H$. Crucially, an exponential numerical sensitivity appears only
\emph{inside a logarithm}. In particular,
\[
 H=O\!\left(L^2K^2r^{-4}(d+\log L)^2
 \log\frac{K(d+\log L)}{r}\right)
\]
up to universal changes of constants and harmless additions inside logarithms.

\section{Learning and nonvacuous circuit targets}

\subsection{Robust digital gates}

Consider a layered Boolean circuit with at most $C$ wires per stage, including
a constant $+1$ rail. Carry unused wires by identity gates and fill unused
positions by constants. AND, OR, NOT and COPY are sign functions of
$a+b-c$, $a+b+c$, $-a$, and $a$, where $c=+1$.
Replace each linear form by its tanh after multiplication by $8$, and use
readout coefficient $2$ on the answer wire.

If every incoming noisy coordinate differs from its Boolean value by at
most $1/8$, each intended signed preactivation is at least $5$.
The tanh output then has error $<1/16$. Adding a further coordinate
perturbation of magnitude at most $1/16$ restores the same invariant.
Thus errors do not accumulate with depth. A base RMS radius
$1/(16\sqrt C)$ suffices, its energy is at most
\[
 B_0=576d+576(L-2)C+4C=O(d+LC),
\]
and its noisy output margin exceeds one. Repeated inputs can give an edge
of magnitude $16$; pad to $C\ge16$ if the bounded-weight convention requires
weights no larger than width.

At every $N\ge C$, clone the wires equally and leave dead padding.
Clone-group means follow the original computation with averaged injections.
The normalized total noise increases by at most $\sqrt2$ under this map,
and the norm increases by at most two. Hence every such learner width has
a target witness of energy at most $2B_0$ and all-layer radius at least
$1/(16\sqrt{2C})$. This is an explicit robust target construction, rather
than an assumption that ordinary analog complexity automatically equals
robust complexity.

\subsection{Strong-prior theorem}\label{sec:strong-limit}

For fixed $r$, let $b_{r,N}(g)$ be the infimum of $\kappa_N$ on
$\mathcal R_{r,N,g}$. These implementation events are open: the noise ball
is compact, the input set is finite, and the minimum noisy signed output
is continuous in the weights. The usual Gaussian small-variance argument
therefore gives
\[
 \lim_{\sigma\downarrow0}2\sigma^2\log
 \mu_{N,\sigma}(\mathcal R_{r,N,g})=-b_{r,N}(g).
\]
The normalization of the robustness-conditioned prior cancels in ratios.
For completeness, this rate calculation is elementary in the fixed
finite-dimensional parameter space. After a linear change of coordinates,
$\kappa_N$ is a squared Euclidean norm. The Gaussian Chernoff bound with
parameter $(1-\xi)/(2\sigma^2)$ gives rate at most
$-(1-\xi)\inf_U\kappa_N$ for any measurable event $U$; let $\xi\downarrow0$.
For an open nonempty $U$, a small ball about a point of energy within $\eta$
of its infimum has positive volume and maximum energy arbitrarily close
to that infimum. Integrating the Gaussian density on that ball gives the
matching lower rate as $\eta\downarrow0$.

\begin{theorem}\label{thm:strong}
Let $0<r\le1$, $K\ge1$, and $\varepsilon,\delta\in(0,1)$.
Suppose $b_{r,N}(f)\le K$. There is a threshold $\sigma_0>0$, depending on
$d,L,N,r,\phi,f,\varepsilon$ but not on the input distribution, sample size,
or $\delta$, such that for every $0<\sigma<\sigma_0$ it suffices to take
\[
 n\ge\frac2\varepsilon\left[H(d,L,K+1,r)+\log\frac1\delta\right]
\]
to obtain posterior average error at most $\varepsilon$, with probability
at least $1-\delta$ over IID training examples from any input distribution.
The sufficiently-small-variance threshold may depend on width.
\end{theorem}

\begin{proof}
There are finitely many Boolean functions. The total prior mass of functions
with $b_{r,N}(g)>b_{r,N}(f)$, divided by the target mass, tends to zero.
The functions with rate no larger than the target have log cardinality at
most $H(d,L,K+1,r)$, by near-minimizers and Theorem~\ref{thm:compression}.
A union bound excludes every member of this finite class with error
$>\varepsilon/2$ at the stated sample size. The remaining posterior mass
outside the class is at most $\varepsilon/2$, since the target is always
consistent. This proves the claim.
\end{proof}

The digital construction makes this bound polynomial jointly in $d,L,C$
at inverse-polynomial $r$, uniformly over $N\ge C$ in the original order of
limits. A fixed radius prior covers a specified class of circuit widths;
it is not silently changed for each unknown target.

\subsection{One fixed prior without a target-size choice}\label{sec:selfcal}

A possible calibrated support condition uses an aggregate of layer norms,
rather than their product. Define
\[
 S_N(w)=\frac dN\|W^1\|_F^2+
 \sum_{s=2}^{L-1}\|W^s\|_\op^2+N\|v\|_2^2\le\kappa_N(w).
\]
Require clean margin $>1/2$ together with all-layer sign robustness at radius
\[
 r(w)=\frac{c}{\sqrt{1+S_N(w)}},\qquad c=1/100.
\]
This is a proposed prior design, rather than an empirically established
normalization. It expresses an inverse-polynomial internal precision
requirement relative to an aggregate measure of layer amplification.
The rule is independent of the target, its size, the data, and $\gamma$.
It charges for internal precision rather than forbidding particular Boolean
functions outright: increasing the readout norm can lower the required
radius, but the Gaussian prior charges for that increase.
It remains open in parameter space: write the perturbation as $r(w)u$
over the fixed compact unit ball, and use continuity of $S_N$.
On the energy sublevel $\kappa_N\le B+1$, every member has radius at least
$c/\sqrt{B+2}$, so the compression and counting bounds are polynomial
in $d,L,B$. The cloned gate witness has $S_N\ge4C$ from its
readout alone, making the required radius smaller than its certified radius.
Its energy is at most $2B_0$. Thus this \emph{single} support rule gives the
strong-prior polynomial-depth conclusion for all layered Boolean circuits,
without a target-complexity-dependent choice of $r$. The same rule also works
uniformly at every $\gamma\ge1$, as proved below. Using $\kappa_N$ instead of
$S_N$ in this calibration would suffice in the strong-prior limit, but would
not give the same finite-variance argument throughout $1<\gamma<2$.

\subsection{Every fixed exponent \texorpdfstring{$\gamma>1$}{gamma > 1}}

\begin{theorem}\label{thm:gamma}
Fix $\gamma>1$. Let $f$ be computed by a layered fan-in-two Boolean circuit
of width $C$ and $L-1$ stages, and let
$0<r\le1/(64\sqrt{2C})$. Condition the specified Gaussian prior on
$\bigcup_g\mathcal R_{r,N,g}$. For every learner width $N\ge C$,
polynomially many samples in
\[
 d,\ L,\ C,\ 1/r,\ 1/\varepsilon,\ \log(1/\delta)
\]
suffice for posterior average error at most $\varepsilon$ with probability
$1-\delta$. The polynomial degree may depend on $\gamma$, but not on $L$.
\end{theorem}

Here are the probability and robustness steps; explicit constants for the
target-mass bound appear in Appendix~\ref{sec:cylinder}.
The robust digital cylinder gives, for every $\gamma\ge1$,
\begin{equation}\label{eq:targetmass}
 \mu_{N,\gamma}(\mathcal R_{r,N,f})\ge e^{-DN^\gamma},
 \qquad D=\poly(d,C,L).
\end{equation}
For a tail parameter $K\ge1$, put $a=r/(4\sqrt L)$.
The Gaussian singular-value-tail argument in
\emph{Width-uniform Bayesian generalization above the mean-field
scaling}~\cite{aboveone} implies that, once
\[
 N^{\gamma-1}\ge\max\{8(d+1),64/a^2\},
\]
discarding singular values at most $a$ from the internal matrices leaves
energy at most $LK$, outside raw probability $Le^{-KN^\gamma/8}$.

Robustness converts the discarded spectral part into bounded activation
noise without a downstream product of operator norms.
Along any radius-$r/2$ noisy pass of the truncated network, hidden states
have norm at most $(1+r/2)\sqrt N$. Each discarded matrix produces an
additional normalized postactivation error at most $a(1+r/2)$. Combining
all layers gives total noise at most
\[
 r/2+\sqrt L\,a(1+r/2)\le7r/8<r.
\]
That exact trajectory is therefore an allowed noisy pass of the original
draw. The truncated network retains all-layer sign margin $r/2$, though
not necessarily clean margin $1/2$. Theorem~\ref{thm:compression} restores
the clean margin and places its sign function in a class with entropy
$H(d,L,LK,r/2)$.

Choose
\[
 K=8[D+\log(2L/\varepsilon)],\qquad
 T=\left\lceil\max\{C,[8(d+1)]^{1/(\gamma-1)},
 (32\sqrt L/r)^{2/(\gamma-1)}\}\right\rceil.
\]
For $N\ge T$, the exceptional mass divided by target mass is at most
$\varepsilon/2$. For $C\le N<T$, the negative logarithm of the target mass is
at most $DT^\gamma$. To use this directly, integrate over the prior:
a fixed function of error greater than $\varepsilon/2$ agrees with all
$n$ labels with probability at most $e^{-n\varepsilon/2}$.
Divide this expected consistent bad mass by the target-mass lower
bound and apply Markov's inequality. Consequently the sufficient bound
\[
 n\ge\frac2\varepsilon\left[
 \max\{H(d,L,LK,r/2),DT^\gamma\}
 +\log\frac{2}{\varepsilon\delta}\right]
\]
proves the theorem. Every displayed term has a polynomial majorant of the
claimed form. These estimates deteriorate as $\gamma\downarrow1$;
Section~\ref{sec:unified} gives a different, uniform argument.

\begin{corollary}\label{cor:singleprior}
The single calibrated support condition of Section~\ref{sec:selfcal}
has polynomial sample complexity jointly in $d,L,C$ at every fixed
$\gamma>1$, for every learner width $N\ge C$, for the same circuit targets.
There is no target-size-dependent robustness parameter or hard norm cutoff.
\end{corollary}

\begin{proof}
The target cylinder from Appendix~\ref{sec:cylinder} still lies in the
calibrated support. Indeed its readout center satisfies
$\sqrt N\|v_*\|\ge\sqrt{2C}$ and the cylinder changes this norm by at most
$1/4$. Therefore $S_N\ge(\sqrt{2C}-1/4)^2\ge C$ throughout the cylinder,
and $r(w)\le c/\sqrt{1+C}<1/(64\sqrt{2C})$. Its raw mass is still
at least $e^{-DN^\gamma}$.

Choose $K=8[D+\log(2L/\varepsilon)]$, put
$r_0=c/\sqrt{1+LK}$, and use cutoff $a=r_0/(4\sqrt L)$.
On the same Gaussian good event as in Theorem~\ref{thm:gamma}, the first
and readout energy terms are at most $K$, and each internal matrix's
retained singular values have squared sum at most $K$.
Every original singular value is then at most $\max\{\sqrt K,a\}=\sqrt K$,
so $S_N\le LK$. Thus every supported good draw has robustness radius
at least $r_0$, and its truncation has energy at most $LK$ and radius
$r_0/2$. Its own calibrated radius need not be retained; the fixed-radius
compression theorem is all that is required.

Use the preceding sample bound with $H=H(d,L,LK,r_0/2)$ and
\[
 T=\left\lceil\max\left\{C,[8(d+1)]^{1/(\gamma-1)},
 \left[\frac{1024L(1+LK)}{c^2}\right]^{1/(\gamma-1)}\right\}\right\rceil.
\]
The target-relative tail and the small-width argument are unchanged.
Every term has the claimed polynomial majorant.
\end{proof}

\subsection{One theorem for all \texorpdfstring{$\gamma\ge1$}{gamma at least one}}
\label{sec:unified}

The preceding spectral truncation proof does not extend uniformly to mean
field: the Gaussian internal bulk has operator norm of order one, and keeping
all singular values above a small cutoff can leave energy of order $N$.
Appendix~\ref{sec:robust-entropy} resolves this by compressing a posterior
distribution, while tracking conditional entropy only in the original weights
that have not yet been processed. All approximation errors become simultaneous
activation perturbations of the original network. The proof allows arbitrary
mixtures of functions in the posterior.

\begin{theorem}[A single robustness restriction for all widths and scales]
\label{thm:unified}
Use $\tanh$ activation and the support restriction of
Section~\ref{sec:selfcal}: global clean margin $1/2$ and all-layer sign
robustness at radius $r(w)=1/[100\sqrt{1+S_N(w)}]$.
Let $f$ be computed by a layered fan-in-two Boolean circuit over
AND, OR, NOT, COPY, and constant gates, of width $C\ge1$ with $L-1$ stages,
including the constant rail. There is a fixed
polynomial $P$ such that, for any input distribution on $\X$, any
$0<\varepsilon,\delta<1$, and
\[
 n\ge P(d,L,C,\varepsilon^{-1},\log(1/\delta)),
\]
with probability at least $1-\delta$ over the noiseless sample, the posterior
Gibbs error is at most $\varepsilon$, simultaneously for every $N\ge C$ and
every raw Gaussian with entry variances $t^2/d,t^2/N,t^2/N^2$, $0<t\le1$.
The same polynomial, after a constant rescaling of $\varepsilon$, bounds the
fraction of inputs where the posterior probability of the correct label is
at most $2/3$, or at most $1/2$.
In particular the result is uniform in all $\gamma\ge1$, by taking
$t^2=N^{1-\gamma}$, and in every strong-prior scale $\sigma^2\le1/N$.
Neither coefficients nor degree of $P$ depend on $N,t$, or $\gamma$.
\end{theorem}
\begin{proof}
Write $\mu=\mu_{N,t}$ for the raw Gaussian. The robust target cylinder of
Appendix~\ref{sec:cylinder} has mass at least $\exp(-DN/t^2)$ for a
$D=\poly(d,L,C)$, uniformly for $N\ge C$, $0<t\le1$. Its radius exceeds
the required calibrated radius, as checked in Section~\ref{sec:selfcal}.
For every sample $s$, the posterior $Q_s$ is $\mu$ conditioned on the support
restriction and consistency with $s$. Hence
\[
 D_{\rm KL}(Q_s\Vert\mu)\le DN/t^2.
\]
Set $\beta=\varepsilon/4$, $\alpha=\varepsilon/12$, and
\[
 B=\max\{16,8(d+1),8[D+\log(L/\beta)]\},\quad
 G=\{S_N\le LB\},\quad r_0=\frac{1}{100\sqrt{1+LB}}.
\]
The first-layer term of $S_N$ is $(t^2/N)\chi^2_{N(d+1)}$ and the readout
term is $(t^2/N)\chi^2_N$. Their Gaussian moment bounds, and
$\Pp(\|W^s\|_{\op}>2t+u)\le\exp[-Nu^2/(2t^2)]$ internally, give
\[
 \mu(G^c)\le L\exp[-NB/(8t^2)],\qquad Q_s(G^c)\le\beta.
\]
Consequently $Q_s^G=Q_s(\cdot\mid G)$ has all-layer radius at least $r_0$,
is still consistent, and has KL at most $N(D+1)/t^2$.

Theorem~\ref{thm:robust-entropy} supplies one finite class
$\mathcal H(d,L,D+1,r_0,\alpha)$, independent of $N,t,s$, with polynomial
log-size $H_0$, and a law $\Gamma_s$ on it. At each input, $\Gamma_s$ and
$Q_s^G$ admit a prediction coupling with mismatch probability at most
$\alpha$. Thus the expected empirical error of $\Gamma_s$ is at most
$\alpha$, and its true error differs from that of $Q_s^G$ by at most $\alpha$.
With probability $1-\delta$, simultaneously for every $h\in\mathcal H$,
\[
 \operatorname{err}(h)\le2\operatorname{err}_s(h)
            +\frac{8[H_0+\log(1/\delta)]}{n}.
\]
Indeed, violation implies empirical error below half the true error and
true error above $8[H_0+\log(1/\delta)]/n$; the multiplicative Chernoff
bound and a union bound apply. Averaging under $\Gamma_s$ gives
\[
 \operatorname{GibbsErr}(Q_s)
 \le\beta+3\alpha+\frac{8[H_0+\log(1/\delta)]}{n}\le\varepsilon
\]
when $n\ge16[H_0+\log(1/\delta)]/\varepsilon$. All parameters of $H_0$
have polynomial bounds in the stated quantities. The same finite-class event
works for every $N,t$. Finally, if $p_s(x)$ denotes posterior probability of
the correct label, then
$\Pp_x\{p_s(x)\le q\}\le\operatorname{GibbsErr}(Q_s)/(1-q)$.
\end{proof}

The same compression theorem gives a polynomial bound for a fixed all-layer
radius $r$ whenever the raw robust target mass is at least $\exp(-NK/t^2)$
with $K$ polynomial in target parameters. This is recorded in
Corollary~\ref{cor:robust-entropy-pac}.

\paragraph{A simpler alternative with a hard energy cutoff.}
One may instead restrict the support further to
$\kappa_N<K$, with $K$ independent of learner width. Its entire function
class has entropy at most $H(d,L,K,r)$. If the target has a strict robust
witness of energy below $K$, it has positive prior mass at every nonzero
variance. For
\[
 n\ge[H(d,L,K,r)+\log(1/\delta)]/\varepsilon,
\]
every consistent supported function has error at most $\varepsilon$ with
probability $1-\delta$. This holds for any Gaussian variance, including
mean field. This is a stronger change of prior, and is unnecessary in
Theorem~\ref{thm:unified}. In particular a raw energy budget
that itself grows with width would not provide the advertised width independence.

\subsection{A normalization with a nonexpansive-network interpretation}
\label{sec:readout-reference}

An alternative support rule has a simpler interpretation. Write $J=L-1$,
$E_v=N\|v\|_2^2$, and require the same global clean margin and all-layer
sign robustness, now at radius
\[
 r_{\rm ref}(w)=\frac{1}{100\sqrt{1+J E_v(w)}}.
\]
The denominator is the total-noise gain bound for a network with
nonexpansive hidden transitions and the same readout. This is a fixed
fraction of its guaranteed noise radius, with no target-size or Gaussian-scale
parameter and no hard energy cutoff.

\begin{proposition}\label{prop:nonexpansive-reference}
If every internal hidden matrix satisfies $\|W^s\|_{\op}\le1$ and the
clean output margin exceeds $1/2$, this support condition holds.
The first matrix need not be nonexpansive.
\end{proposition}
\begin{proof}
Let $a_s$ be the Euclidean hidden displacement. There is no input noise,
so $a_1\le\|\eta^1\|_2$; subsequently,
$a_s\le a_{s-1}+\|\eta^s\|_2$ by nonexpansiveness. Thus
\[
 |z(x;\eta)-z(x)|
 \le\sqrt{E_v}\sum_{s=1}^J\frac{\|\eta^s\|_2}{\sqrt N}+|\eta^L|
 \le\sqrt{1+JE_v}\,\|\eta\|_*.
\]
At the specified radius this is at most $.01<1/2$.
\end{proof}

The support itself permits larger operator norms whenever the actual
finite-radius robustness suffices. Hence this is an interpretation and a
sufficient architectural condition, not an additional spectral restriction.

\begin{corollary}\label{cor:readout-reference}
All conclusions of Theorem~\ref{thm:unified} also hold for this
readout-reference support: one polynomial sample bound simultaneously for
every $N\ge C$ and every $0<t\le1$, including all $\gamma\ge1$.
\end{corollary}
\begin{proof}
Every draw in the target cylinder of Appendix~\ref{sec:cylinder} has
$E_v\ge C$ and sign robustness at radius $\tau_0=1/(64\sqrt{2C})$.
Since $J\ge1$,
\[
 r_{\rm ref}(w)\le .01/\sqrt{1+C}<\tau_0.
\]
The same target mass lower bound $\exp(-DN/t^2)$ therefore applies.
For $0<\beta\le1/2$, put
\[
 B=\max\{8,8[D+\log(1/\beta)]\},\qquad G=\{E_v\le B\}.
\]
Under the raw Gaussian, $E_v=(t^2/N)\chi_N^2$. Its exponential moment gives
\[
 \mu_{N,t}(G^c)\le\exp[-NB/(8t^2)],\qquad Q_s(G^c)\le\beta,
\]
where the second bound divides by the target cylinder mass.
The consistent conditional law $Q_s(\cdot\mid G)$ has radius at least
$r_0=.01/\sqrt{1+JB}$ and KL at most $N(D+1)/t^2$.
Theorem~\ref{thm:robust-entropy} and the same finite-class argument as in
Theorem~\ref{thm:unified} give
\[
 \operatorname{GibbsErr}(Q_s)
 \le\beta+3\alpha+8[\log|\mathcal H|+\log(1/\delta)]/n,
\]
with one class $\mathcal H(d,L,D+1,r_0,\alpha)$ independent of $N,t$.
The earlier choices of $\alpha,\beta,n$ prove the claim.
\end{proof}

This normalization is less tied to the intermediate spectral estimates than
$S_N$. The two required radii have no universal ordering: that ordering depends
on whether $S_N$ is larger or smaller than $JE_v$. Both are proposed prior
designs, not established laws of ordinary training. In our saved toy models,
the readout-reference radius is always smaller than the spectral-calibrated
radius, so every successful certificate reported below transfers to it.


\section{Naturalness and what trained networks are known to satisfy}

There is a direct numerical interpretation. Rounding every hidden activation
and the scalar output with absolute error at most $2^{-b}$ creates total
normalized noise at most $\sqrt L\,2^{-b}$. An all-layer margin $r$ therefore
preserves every label whenever
\[
 b\ge\left\lceil\log_2(\sqrt L/r)\right\rceil.
\]
An inverse-polynomial margin means that only logarithmically many fractional
bits of activation precision are needed. This is a consequence of robust
computation; choosing a quantized activation rule alone does not establish
the adversarial finite-radius property.

The most defensible condition is finite all-layer robustness, rather than
a product of raw weight norms. A sufficient, more diagnostic condition is
bounded internal amplification throughout a perturbation tube. For smooth
activations define
\[
 G_r(w)=\sup_{x,\|\eta\|_*\le r}
 \left(1+N\sum_{s<L}\|\nabla_{\eta^s}z_w(x;\eta)\|_2^2\right)^{1/2}.
\]
Integration along $t\eta$ gives output change at most $G_r(w)\|\eta\|_*$.
Thus clean margin $>1/2$ and $rG_r(w)\le1/3$ imply all-layer sign margin $r$.
This uses actual downstream Jacobians, including cancellation and saturation,
and must hold in a finite neighborhood, not only at the clean point.

If transition Jacobian norms on that tube are bounded by $q_s$, then
\[
 G_r(w)^2\le1+N\|v\|_2^2
 \sum_{s=1}^{L-1}\prod_{j=s+1}^{L-1}q_j^2.
\]
Nonexpansive hidden transitions give $G_r\le\sqrt{1+(L-1)\kappa_N}$.
More generally, $q_j\le1+a_j$, $a_j\ge0$, and
$\sum_j a_j=O(\log L)$ give polynomial amplification. These are transparent
sufficient restrictions, but raw spectral nonexpansiveness may exclude
useful circuit implementations. Saturation can supply much better local
certificates. Residual step-size bounds give related architectural examples,
but residual networks are not silently substituted into the original tanh model.

Relevant primary literature includes:
\begin{itemize}
\item Wei--Ma's all-layer margin~\cite{weima2020}: a closely related
perturbation notion, generalization theory, and gradient-based margin
optimization on practical image models.
\item Wei--Ma's Lipschitz augmentation~\cite{weima2019}: polynomial explicit
depth dependence in measured hidden-state and interlayer-Jacobian quantities,
including smooth activations. The remaining matrix norms need not be
width-independent under our $\kappa_N$.
\item Arora--Ge--Neyshabur--Zhang~\cite{arora}: noise stability and compression
in trained networks. Their random/compression-specific and example-local
properties are weaker than our all-input adversarial condition.
\item Wei--Chen--Ma~\cite{weichenma}: robust circuit simulations already play
a role in polynomial circuit-learning bounds for a different learner and
parameter-margin objective. Our Gaussian-prior and width normalization need
their own arguments.
\end{itemize}

These sources support naturalness and trainability. They do not prove that
ordinary gradient descent yields an inverse-polynomial worst-input,
simultaneous RMS robustness radius. The distinction is material:
IID noise and worst-case noise differ by a factor $\sqrt N$ in linearized
response, and an all-input guarantee is much stronger than a typical-data
guarantee. The numerical scripts referenced by Fable's note~\cite{fable}
were unavailable to the 17 October investigation; its reported experiments
are not treated as independently reproduced here.

\subsection{A small reproducible training probe}

We trained 24 predetermined width-32 tanh networks on all 64 inputs of the
six-bit cube: majority and parity on the first three bits, hidden depths
$2,4,8,16$, and three seeds. The optimizer was full-batch Adam, for 6000 steps;
there were no later-layer biases. All fits succeeded, with minimum clean
margins between $0.99607$ and $0.99994$. No failed fit or unusual seed was
removed. This is a toy test of actual gradient-trained networks, not evidence
about large practical models.

For the total noise norm~\eqref{eq:noisenorm}, the following are medians over
seeds of the worst radius across the entire cube. A sufficient radius is
obtained from a nonlinear finite-neighborhood derivative bound; an attack
radius is one at which an analytic perturbation or PGD gives a sign failure.
\begin{center}
\begin{tabular}{llrr}
\toprule
Task & Hidden depth & Sufficient radius & Found attack radius\\
\midrule
Majority$_3$ & 2 & .0751 & .115\\
 & 4 & .0361 & .0790\\
 & 8 & .0156 & .0563\\
 & 16 & .00644 & .0566\\
Parity$_3$ & 2 & .0576 & .0938\\
 & 4 & .0278 & .0741\\
 & 8 & .0119 & .0597\\
 & 16 & .00297 & .0449\\
\bottomrule
\end{tabular}
\end{center}

The analytic sufficient bound is valid, but numerical evaluation uses
float64 SVD rather than interval-certified arithmetic. An unsuccessful
attack is never a robustness certificate. One retained depth-16 parity
seed has spectral product $7.38\times10^6$, sufficient local radius
$4.99\times10^{-5}$, and a found attack at $.0933$; this is a large gap
between available bounds, not a proof of robustness at $.0933$.
The tested depth range and gaps do not establish polynomial scaling.
These numbers record the 17 October experiment. The original investigation
retained scripts, all-seed tables, saved weights, and raw data; this
standalone manuscript package does not include that full experimental
archive or constitute a new reproduction of it.

A separate, subsequently predeclared stress cohort used hidden depths $32,64$
with the same optimizer settings and three seeds for both tasks. All six
depth-32 fits succeeded; only three of six depth-64 fits did. The three
retained failures made $23$, $31$, and $30$ errors on the 64 training inputs,
so their robustness radius relative to the target labels is zero; this does
not assert that robustness of their own clean predictions is zero.
Among the successful depth-64
fits, sufficient sign radii ranged from $1.10\times10^{-5}$ to
$4.45\times10^{-4}$, while known attack radii ranged from $.01195$ to $.120$.
This cohort is reported separately, not pooled with the original grid.
It reinforces the limitation: neither optimization success nor an
inverse-polynomial robustness law follows from ordinary training alone.

We also tested the exact support condition used in
Theorem~\ref{thm:unified}, recomputing $S_N$ from every saved model and
evaluating the analytic bound at $r=.01/\sqrt{1+S_N}$.
Of the original 24 models, 23 pass this sufficient test and one is unresolved.
In the 12-model stress cohort, five pass, three violate the clean absolute
margin requirement, and four are unresolved by that test. A stronger
finite-radius certificate, derived in Appendix~\ref{sec:clean-transport-certificate},
retains exact clean Jacobian products and bounds nonlinear Taylor remainders.
Two independent implementations pass all five previously unresolved models,
including the three successful depth-64 fits. Thus all 33 successful fits
pass a sufficient test for the exact proposed support, and the three failed
fits violate its clean absolute margin, subject throughout to the
floating-point qualification. The earlier and improved audits are both
retained. These tests preserve each network's own clean labels, as required
by the prior; a training error alone is not used as a robustness counterexample.
For these 33 fits, the readout-reference radii imply the conservative
activation-rounding guarantee above with $11$--$15$ fractional bits,
subject to the same numerical qualification.

\section{What the necessity question still requires}

The existing norm construction represents every truth table on $q$ bits with
\[
 \kappa_N\le A\left[dq^2\log^2(2q)+(L-2)2^{2q/(L-2)}\right],
 \qquad N\ge2^{q+1}.
\]
At $q=d$, $L=d+2$, this is polynomial norm but $2^{2^d}$ functions.
Appendix~\ref{sec:inputcounter} strengthens it to global polynomial input
Lipschitz constants. This invalidates a polynomial joint norm-capacity bound.

It does not establish the desired Gaussian posterior lower bound. Most
truth tables have large target complexity. A large low-norm class does not
show that the relevant competitors have lower rate than a particular simple
target, nor that their integrated Gaussian mass is large enough after data.
In particular, the common cost of transmitting a signal through many layers
may cancel in target-relative rate comparisons.

There is also a positive statement for the \emph{original} Gaussian prior.
Appendix~\ref{sec:simpletargets} proves that its globally minimum-energy
Boolean functions are exactly $\{\pm1,\pm x_1,\ldots,\pm x_d\}$, at every
depth and width. The strong-prior function prior converges to the uniform
distribution on these $2(d+1)$ functions. Consequently each of these targets
has sample bound $O((\log d+\log(1/\delta))/\varepsilon)$ independent of
depth and width; for uniform inputs, $O(\log(d/\delta))$ samples identify
the target exactly in the limit. The sufficiently-small-variance threshold
can depend on depth and width. This does not settle general circuit targets,
but rules out using the constant target as a strong-prior Gaussian
counterexample merely by invoking the huge low-norm class.
The common minimum rate is also quantitatively understood: writing
$D_0=\max\{d,4\}$, it lies between $L(d/4)^{1/L}$ and
$(L+12)(D_0/4)^{1/L}$. Thus it is $L+O_d(1)$ at large depth, even though
the corresponding sample bound is independent of depth.

Appendix~\ref{app:excessenergy} further separates energy, robustness, and
statistical difficulty. Energy $L+B$ bounds the total squared error of best
rank-one approximations to the internal matrices by $B-\log(d/4)$.
Yet a constant-function network can have energy at most $K_*(4,L)+76$
and be erased by hidden RMS noise as small as $L^{-1}e^{-\sqrt L/4}$.
Thus bounded excess energy does not imply polynomial robustness of every
cheap realization. Conversely, appending layers efficiently shows that
every fixed representable function has infimum energy $L+O(1)$ as depth
grows, with the offset depending on the function, input dimension, and width.
Exponential growth of minimum target energy cannot itself explain a
depth-dependent statistical lower bound.

The distinction can be made exact. A \emph{different} prior assigning equal
mass to all truth tables in the construction needs $\Omega(2^d)$ examples
even for the constant target: every unseen input has posterior probability
$1/2$, so its $>2/3$ confidence-failure fraction is at least $1-n/2^d$.
Small parameter neighborhoods can replace the atoms while preserving their
cube labels and polynomial energy bounds, without changing this statistical
conclusion. This does not assert stability of the global input-Lipschitz
bound under arbitrary parameter perturbations. It is a norm-supported-prior
counterexample, not a counterexample
to the centered independent Gaussian prior.

Conversely, a circuit-description prior mixed with a positive full-support
Gaussian component can retain fragile networks in support and still give
polynomial-depth sample bounds through target prior mass. Thus hard robustness
support is not logically necessary for Bayesian learning in general.
The question addressed by this report concerns the particular Gaussian
learner; the conditional obstruction in~\cite{conditional2026} is a later
result about that same learner.

The reported experiments leave the scaling of finite-radius certificates
with depth and width undetermined. The mathematical results leave an
unconditional unrestricted-Gaussian classification undetermined as well;
the later conditional obstruction does not supply an explicit exponential
simple-target lower bound. Replacing all-input robustness by typical-input
robustness also needs an additional statistical argument, since an
unconstrained exceptional set can contain the training data.

\appendix
\section{Robust Gaussian target mass}\label{sec:cylinder}

This appendix supplies the denominator used in Theorems~\ref{thm:gamma}
and~\ref{thm:unified}, for every $N\ge C$ and every $0<t\le1$.
Take balanced clone groups of sizes $k_a$, so $N/k_a\le2C$, and let $P$
project onto group-constant vectors. Every unused circuit slot is filled by
a Boolean constant, so the ideal wire vector $b^s(x)$ lies in the range of
$P$ and has norm $\sqrt N$.
Put
\[
 e_*=(16\sqrt{2C})^{-1},\quad \tau_0=e_*/4,\quad
 b_L=\sqrt{16\log9+8\log(4L)},
\]
\[
 A=\max\{32b_L,2\log(8/e_*)\},\qquad
 \rho=\min\{1/4,A/[4(d+1)],Ae_*/(64\sqrt C)\}.
\]
Use tanh gate gain $A$ and readout coefficient $2$. The gain need not be
bounded by logical circuit width; it can be accommodated by polynomial
padding if invoking the ``weights bounded by width'' complexity convention.

Allow any $0<t\le1$. The standardized coordinates
$\sqrt d W^1$, $\sqrt N W^s$, and $Nv$ have common standard deviation $t$;
the exponent-$\gamma$ case is $t=N^{(1-\gamma)/2}$.
Restrict every first-layer coordinate, every
internal row's projection onto each normalized clone-group direction, and
every readout coordinate to an interval of radius $\rho t$ around its clone
value. Then
\[
 W^s=W_*^s+R^s+E^s,\qquad R^s=R^sP,
 \quad\|R^s\|_\op\le\rho t\sqrt C,\quad E^sP=0.
\]
The free Gaussian bulk $E^s$ is independent of the cylinder. A $1/4$-net
Gaussian operator bound makes $\|E^s\|_\op\le b_Lt$ simultaneously with
conditional probability at least $1/2$.

We check robustness to a separate per-layer RMS budget $\tau_0$, which is
stronger than the total-budget requirement for $r\le\tau_0$. Write
$e_s=\|h^s_\eta-b^s\|_2/\sqrt N$. If $e_{s-1}\le e_*$, each clone-group
mean has error at most $1/16$, so the teacher part of every intended signed
preactivation is at least $A/2$. For every real $u$ and $\epsilon\in\{-1,1\}$,
\[
 \epsilon a\ge A/2\quad\Longrightarrow\quad
 |\tanh(a+u)-\epsilon|\le2e^{-A/2}+8|u|/A.
\]
This is a global secant bound: for $|u|<A/4$ use saturation, and otherwise
use the trivial bound $2$. It handles concentrated adversarial noises.
Since $E^sb^{s-1}=0$ and noisy states have norm at most
$(1+\tau_0)\sqrt N$,
\[
 e_s\le2e^{-A/2}+(8/A)\rho\sqrt C(1+\tau_0)
 +(8b_L/A)e_{s-1}+\tau_0\le e_*.
\]
The first layer starts the invariant because its coordinatewise
preactivation error is at most $\rho(d+1)\le A/4$.
The readout has noisy signed value at least
$2-1/8-\rho(1+\tau_0)-\tau_0>1$.

There are $Nq$ constrained independent scalar coordinates, where
$q=d+2+(L-2)C$. Their squared centers sum to at most $NB_0'$, with
\[
 B_0'=9dA^2+18(L-2)CA^2+8C.
\]
The elementary Gaussian interval bound and the independent operator event give
\[
 -\log\mu_{N,t}(\mathcal R_{r,N,f})
 \le \frac{N}{t^2}\left[B_0'+q\log\frac{e\sqrt{2\pi}}{2\rho}+\log2\right]
 =DN/t^2.
\]
Here $D$ is jointly polynomially bounded in $d,C,L$. Exponentially small
coordinate precision would in any case cost only its logarithm in this
negative log-mass estimate; it should not be replaced unnecessarily by its
reciprocal.

% Standalone theorem/proof fragment. Uses amsmath, amssymb, amsthm.
\section{Input robustness does not control the depth-dependent norm capacity}
\label{sec:inputcounter}

The following statement concerns norm-class capacity, not the posterior of the
independent Gaussian weight prior.

\begin{proposition}
For every integer $d\ge2$, put $N=2^{d+1}$ and $L=d+3$.
Every $f:\{-1,1\}^d\to\{-1,1\}$ has a width-$N$, depth-$L$ tanh network,
with a linear readout and no biases except the constant input coordinate,
whose real-valued output $z:\mathbb R^d\to\mathbb R$ satisfies
\[
 z(x)=f(x)\quad(x\in\{-1,1\}^d),\qquad
 \kappa_N(z)\le100d^3\log^2(16d),\qquad
 \operatorname{Lip}(z)\le300d^{3/2}\log^2(16d).
\]
Here $\kappa_N$ is the weighted Gaussian energy defined in this report's model.
In particular the correct signed output remains at least $3/4$ under every
Euclidean input perturbation of radius
$[1200d^{3/2}\log^2(16d)]^{-1}$ around a Boolean vertex.
Nevertheless this witness fails the hidden-activation robustness condition
at RMS radius
\[
 \tau_d=\frac{3\sqrt2}{4}\,2^{-d}.
\]
Thus polynomial bounds on weighted norm and even global input Lipschitz
constant do not give a norm-class counting bound polynomial jointly in depth
and the displayed parameters.
\end{proposition}

\begin{proof}
Put $M=2^d$, $m=d-1$, $s=3/4$, $a=\tfrac12\log(16d)$,
$B=\log(4d)$, $c=\tanh B$, and $b=\tanh1$.
For $p\in\{-1,1\}^d$ define
\[
 T_p(u)=\frac{1+\tanh(a(p\cdot u-(d-1)))}2.
\]
The $M\times M$ matrix $T=(T_p(x))_{x,p}$, evaluated on the Boolean cube,
satisfies
\[
 \|T-I\|_\infty
 \le e^{-2a}+e^{2a}\bigl[(1+e^{-4a})^d-1\bigr]<\frac14.
\]
Indeed the first term bounds its diagonal defect and the second sums the
positive off-diagonal terms by Hamming distance.
Consequently $v=T^{-1}f$ satisfies $\|v\|_\infty\le4/3$.
Write $F(u)=\sum_pv_pT_p(u)$.

The first hidden layer contains $r=\lfloor N/(d+1)\rfloor$ copies of
$\tanh(Bx_i)$ for each $i=1,\ldots,d$, and $r$ copies of the constant $c$;
any unused neurons are zero. The next layer has $M$ detector neurons. For
detector $p$, give each copy of coordinate $i$ weight $ap_i/(rc)$ and each
constant copy weight $-a(d-1)/(rc)$. This computes
$\tanh(a(p\cdot u(x)-(d-1)))$, where
\[
 u_i(x)=\frac{\tanh(Bx_i)}c.
\]
Its other $M$ neurons compute the constant $b$, using weight $1/(rc)$ from
each of the $r$ constant predecessors.

Set
\[
 t_0=s2^{-(d-2)/2},\qquad
 R=2^{(d-2)/(2(d-1))},\qquad t_j=t_0R^j\quad(0\le j\le m).
\]
The third hidden layer has one active neuron with weights
$a_p=\tfrac12\operatorname{atanh}(t_0)v_p$ from the detectors and common
weight $(\sum_pa_p)/(Mb)$ from the constant neurons. It computes
$\tanh(\operatorname{atanh}(t_0)F(u(x)))$.
For each of the next $m$ hidden layers put $n_j=\lceil N^{j/m}\rceil$ active
neurons, where $n_0=1$. Each has incoming weight
$\operatorname{atanh}(t_j)/(n_{j-1}t_{j-1})$ from every active predecessor.
Finally use readout weight $1/(Ns)$ on each last-layer neuron.
Since $u$ fixes the Boolean vertices and $F(x)=f(x)$ there, induction shows
that the active neurons of stage $j$ equal $t_jf(x)$ and the output is $f(x)$.
There are $3+m=d+2$ hidden layers, hence depth $L=d+3$.

For completeness, the norm accounting has no hidden exponential factor.
The first layer costs at most $dB^2$. Since $r\ge N/[2(d+1)]$, the second
layer costs at most
\[
 \frac{d+1}{c^2}\bigl[a^2(d+(d-1)^2)+1\bigr]
 \le2(d+1)\bigl[a^2(d^2-d+1)+1\bigr].
\]
The third-layer cost is at most
$(1+b^{-2})(\operatorname{atanh}(s)/s)^2<5$, using $Mt_0^2=9/4$.
Each amplification layer costs at most
\[
 2\Bigl(\frac{\operatorname{atanh}(s)}s\Bigr)^2N^{1/m}R^2
 =2\Bigl(\frac{\operatorname{atanh}(s)}s\Bigr)^2
     2^{2+1/(d-1)}<27,
\]
and the readout costs $s^{-2}=16/9$. These bounds imply the stated
$100d^3\log^2(16d)$ bound with ample slack.

We next bound input sensitivity on all of $\mathbb R^d$. One has
$c^{-1}\le1+1/(7d^2)$ and $\operatorname{Lip}(u)\le B/c<2B$.
For $|u_i|\le c^{-1}$,
\[
 \begin{split}
 \sum_p e^{2a(p\cdot u-(d-1))}
 &\le16d\exp\!\left(\frac1{256d}+2ad(c^{-1}-1)\right)\\
 &\le16d\exp\!\left(\frac1{256d}+\frac{\log(16d)}{7d}\right)<21d.
 \end{split}
\]
Using $\operatorname{sech}^2(t)\le4e^{2t}$ and $\|v\|_\infty\le4/3$ gives
$\|\nabla F(u)\|_2\le56ad^{3/2}$ throughout that box.

The remainder of the network is a scalar map $H$, so $z=H\circ F\circ u$.
By $|\tanh'|\le1$ and telescoping the scale factors,
\[
 \operatorname{Lip}(H)
 \le\prod_{j=0}^{m}\frac{\operatorname{atanh}(t_j)}{t_j}.
\]
For $0<t\le s$, the power series gives
$\log(\operatorname{atanh}(t)/t)\le t^2/[3(1-s^2)]=16t^2/21$.
If $d\ge3$, then $R^2\ge\sqrt2$ and
$\sum_jt_j^2\le s^2/(1-1/\sqrt2)$, hence
\[
 \operatorname{Lip}(H)
 \le\exp\!\left(\frac3{7(1-1/\sqrt2)}\right)<5.
\]
The case $d=2$ has two values $t_j=s$ and satisfies the same bound.
Multiplying the three Lipschitz bounds yields
\[
 \operatorname{Lip}(z)\le560aBd^{3/2}
 \le280d^{3/2}\log^2(16d)<300d^{3/2}\log^2(16d).
\]
The claimed input-perturbation guarantee follows from $z(x)=f(x)$.

Finally, at any Boolean input the third hidden layer has just one nonzero
coordinate, $t_0f(x)$. An additive post-activation perturbation
$-t_0f(x)$ in that coordinate sets the layer to zero. All subsequent hidden
activations and the output are then zero. Its Euclidean norm is $t_0$, and
\[
 \frac{t_0}{\sqrt N}=\frac{3\sqrt2}{4}2^{-d}.
\]
Thus the internal robustness condition fails at that RMS radius.
The $2^{2^d}$ truth tables give the counting conclusion.
\end{proof}

\paragraph{Scope.}
This rules out polynomial-depth \emph{norm-class counting} under polynomial
input-Lipschitz bounds, and gives an explicit internal-fragility mechanism.
It does not prove a sample lower bound for a simple target under the original
independent Gaussian parameter prior: that would require a posterior mass or
minimum-energy comparison specific to that prior.

% Standalone theorem/proof fragment. Uses amsmath, amssymb, amsthm.
\section{Polynomial-overhead robustification}
\label{sec:digitize}

\begin{proposition}
Suppose a width-$C$, depth-$L$ tanh network, with every weight bounded in
absolute value by $C$, implements $f:\{-1,1\}^d\to\{-1,1\}$ with strict
output margin $1/2$. Then $f$ has a tanh implementation of width and depth
polynomial in $(d,C,L)$ that tolerates simultaneous adversarial hidden-layer
activation perturbations of inverse-polynomial RMS size.
More precisely, there is a fan-in-two Boolean circuit for $f$ of size
$S=\operatorname{poly}(d,C,L)$. Put
\[
 W=\max\{8,d+S+1\},\qquad D=S+2.
\]
There is a width-$W$, depth-$D$ tanh network, all of whose weights have absolute
value at most $8$, such that every sequence of post-activation injections with
\[
 \|\zeta^s\|_2\le\frac18\quad(1\le s<D),\qquad
 |\zeta^{\rm out}|\le\frac1{8\sqrt W}
\]
preserves signed output margin strictly greater than $5/8$, on every Boolean
input. Its weighted norm is at most
\[
 \frac{64d(d+1)}W+192SW+W.
\]
At every width $N\ge W$, cloning preserves this guarantee at RMS radius
$1/(8\sqrt{2W})$ and increases weighted norm by at most a factor of two.
The construction allows polynomial depth and width overhead; it makes no
claim at the original depth $L$ or every original width $N\ge C$.
\end{proposition}

\begin{proof}
We first replace the analog computation by a finite Boolean computation.
Round every parameter to absolute precision $\eta=2^{-p}$. Evaluate each
hidden tanh to precision $\eta$ and clip its approximation into $[-1,1]$;
evaluate the readout to the same precision. If $E_s$ is the maximum coordinate
error after layer $s$, the first hidden layer satisfies
$E_1\le(d+2)\eta$, and the subsequent layers and readout satisfy
\[
 E_s\le C^2E_{s-1}+(C+1)\eta.
\]
This follows by writing each affine error as the original weights times
activation errors, plus weight errors times bounded approximate activations.
The clipping operation does not increase the approximation error.
It suffices to choose
\[
 2^{-p}\,[d+2+L(C+1)](C^2+1)^{L-1}\le\frac18.
\]
Thus $p=O(L\log(C+1)+\log(d+L+1))$ bits suffice, and the final sign is unchanged.
The rounded parameters are hardwired constants; no efficient description of
the original real numbers is required.

Here is an elementary polynomial-cost tanh evaluation algorithm. For a dyadic
argument $t$ with $|t|\ge p+2$, return its sign; the error is less than $2^{-p}$.
Otherwise set $x=2|t|\le2(p+2)$ and approximate $e^{-x}$ by its degree
$m=16(p+2)$ Taylor polynomial. Its error is at most
\[
 \frac{x^{m+1}}{(m+1)!}
 \le\left(\frac{ex}{m+1}\right)^{m+1}
 \le(e/8)^{m+1}<2^{-p-4}.
\]
The Taylor sum can be evaluated exactly using the common denominator
$2^{bm}m!$, where $b$ is the number of fractional bits of $x$.
All integers involved have polynomial bit length, even though some Taylor
terms are large and cancel. Round and clip the answer into $[0,1]$, then use
$(1-u)/(1+u)$ and oddness to obtain $\tanh t$.
This rational map has derivative of magnitude at most $2$ on $[0,1]$, so a
constant number of guard bits gives the required accuracy. Integer addition,
multiplication, comparison, and division have polynomial-size Boolean
circuits. Applying this algorithm throughout the rounded network gives a
Boolean circuit of size $S=\operatorname{poly}(d,C,L)$ for $f$.

We now simulate that circuit by noisy tanh gates. Express bits as $\pm1$, and
carry a constant $+1$ channel. Use
\[
 \begin{array}{ll}
 \mathrm{COPY}(a)=\tanh(8a),&\mathrm{NOT}(a)=\tanh(-8a),\\
 \mathrm{AND}(a,b)=\tanh(8(a+b-c)),&
 \mathrm{OR}(a,b)=\tanh(8(a+b+c)),
 \end{array}
\]
where $c$ is the constant channel, itself updated by $\tanh(8c)$.
Simplify gates with repeated or constant inputs beforehand, so each displayed
nonzero edge has magnitude $8$.
If all incoming signals are within $1/4$ of their ideal values, the signed
affine expression for AND and OR is at least $1-3/4=1/4$; COPY and NOT have
larger margin. Consequently each gate output is within
$1-\tanh2<1/8$ of its ideal value. Adding any coordinate perturbation of
magnitude at most $1/8$ leaves the error strictly below $1/4$.
This invariant holds at arbitrarily many layers, with no accumulated error.

Use one Boolean gate per hidden layer and copy all previously computed
signals forward. The first hidden layer initializes the input signals and
constant from $\tanh(8x_i)$, including $x_0=1$.
The $W$ channels accommodate all inputs, gate outputs, and the constant.
Unused channels have zero outgoing weights until used.
Each hidden-layer Euclidean noise budget $1/8$ bounds every coordinate
injection by $1/8$, so the invariant applies. Read out the desired final
signal with coefficient one. It has signed value greater than $3/4$; the
output injection is at most $1/8$, leaving signed margin greater than $5/8$.
The first layer has $d+1$ nonzero edges, each of the $S$ gate layers has at
most $3W$ nonzero edges, and the readout has one edge. This gives the stated
weighted norm.

Finally, at width $N\ge W$ use $k=\lfloor N/W\rfloor$ copies of each hidden
neuron, divide outgoing weights by $k$, and leave the remaining channels dead.
Clone-group means obey exactly the original computation with averaged
post-activation injections. Their Euclidean noise norm is at most
\[
 \frac1{8\sqrt{2W}}\sqrt{\frac Nk}\le\frac18,
\]
because $N/k\le2W$. Thus the same invariant and margin hold.
The first-layer norm does not increase, internal-layer norm is unchanged,
and the weighted readout norm increases by at most two.
\end{proof}

\paragraph{Interpretation.}
Inverse-polynomial internal robustness does not prevent representation of
any Boolean function already having a bounded-weight tanh implementation,
provided polynomial increases in depth and width are allowed. This is an
expressivity statement, not a claim that gradient descent produces the
constructed digital simulation. The point is that exponentially small analog
precision costs only polynomially many digital bits.

% Standalone theorem/proof fragment. Uses amsmath, amssymb, amsthm.
\section{A depth-independent Gaussian result for the simplest targets}
\label{sec:simpletargets}

The norm-capacity obstruction does not automatically make even the simplest
targets hard for the specified Gaussian prior. In fact constants and dictators
have a depth-independent sample bound in the strong-prior limit.

\begin{theorem}
Fix $d\ge1$, $N\ge1$, and $L\ge2$, with tanh hidden activations, the constant
input coordinate, and a linear readout as defined in this report's model. Let
\[
 \mathcal M_d=\{+1,-1,+x_1,-x_1,\ldots,+x_d,-x_d\}.
\]
Among all Boolean functions implemented with the all-input strict margin
$1/2$, the functions of smallest weighted norm are exactly $\mathcal M_d$.
Their common infimum norm value depends on $d,L$ but not on $N$.

Consequently, for every $f\in\mathcal M_d$ and every $\varepsilon,\delta\in(0,1)$,
there is $\sigma_0>0$ such that the original globally margin-conditioned
Gaussian posterior has posterior-draw risk at most $\varepsilon$, with
probability at least $1-\delta$, whenever
\[
 0<\sigma<\sigma_0,\qquad
 n\ge\frac2\varepsilon
       \left(\log(2(d+1))+\log\frac1\delta\right).
\]
The guarantee holds for arbitrary input distributions on the Boolean cube.
The sample bound is independent of depth and learner width.
The variance threshold may depend on $d,L,N,\varepsilon,f$; no uniform
threshold over depths or widths is asserted.
\end{theorem}

\begin{proof}
First, the function $Q(u)=\tanh^2\sqrt u$ is strictly increasing and strictly
concave for $u\ge0$. To see concavity, put $t=\sqrt u>0$ and $y=\tanh t$.
A direct derivative computation gives
\[
 Q''(u)=\frac{1-y^2}{2t^3}\bigl[t(1-3y^2)-y\bigr]<0.
\]
Indeed $y-t(1-y^2)>0$ for $t>0$, since it vanishes at zero and has derivative
$2ty(1-y^2)>0$; hence the bracket is less than $-2ty^2$.
The continuous extension at zero preserves concavity.
Jensen's inequality therefore implies, for every $v\in\mathbb R^N$,
\begin{equation}\label{eq:tanh-rms}
 \frac{\|\tanh v\|_2}{\sqrt N}
 \le\tanh\!\left(\frac{\|v\|_2}{\sqrt N}\right).
\end{equation}

For any network weights define
\[
 a=\frac{\|W^1\|_F}{\sqrt N},\qquad
 b_s=\|W^s\|_F\ (2\le s<L),\qquad
 c=\sqrt N\|v\|_2.
\]
Their weighted squared norm is
\[
 \kappa_N=d a^2+\sum_{s=2}^{L-1}b_s^2+c^2.
\]
Define the scalar recursion $r_1=\tanh a$, $r_s=\tanh(b_s r_{s-1})$ for
$2\le s<L$, and let $F(a,b,c)=c r_{L-1}$.
Let $K_*(d,L)$ be the minimum of the displayed cost over nonnegative
$(a,b_2,\ldots,b_{L-1},c)$ with $F\ge1/2$.
The minimum exists: the constraint is closed and nonempty and the cost is
coercive. Every feasible tuple has all entries positive. At a minimizer,
$F=1/2$, since otherwise decreasing $c$ lowers the cost.

Every tuple is realized by a constant-channel network: all first-layer rows
have weight $a$ on the constant input and zero elsewhere; every entry of
$W^s$ is $b_s/N$; and every readout entry is $c/N$.
Its output is the constant $F$ and its norm is the stated cost, at every width.
Moving the first-layer weight from the constant input to input $x_i$ gives
output $F x_i$ by oddness of tanh; negating the readout gives the negative
functions. Thus all members of $\mathcal M_d$ have infimum norm at most $K_*$,
with strict margins obtained by an arbitrarily small increase of the readout.

Conversely, consider a network with $|z_w(x)|\ge1/2$ at every Boolean input.
For uniform $x$ on the Boolean cube and $\bar x=(1,x)$,
\[
 \mathbb E_x\frac{\|W^1\bar x\|_2^2}{N}
 =\frac{\|W^1\|_F^2}{N}=a^2.
\]
There is therefore an $x_*$ with
$\|W^1\bar x_*\|_2/\sqrt N\le a$.
Apply \eqref{eq:tanh-rms} at each hidden layer, the inequality
$\|Wh\|_2\le\|W\|_F\|h\|_2$ at each internal affine map, and Cauchy--Schwarz
at the readout. These give
\[
 \frac12\le|z_w(x_*)|\le F(a,b,c).
\]
Thus its norm is at least $K_*$. In particular, the members of
$\mathcal M_d$ have exactly the global infimum norm $K_*$.

We check the equality case to exclude all other Boolean functions.
Suppose a globally margin-$1/2$ network has norm exactly $K_*$.
Its associated scalar tuple minimizes the scalar problem, so $F=1/2$.
All parameters in that tuple are positive, so the scalar recursion is strictly
increasing in the first preactivation RMS.
Choose $x_*$ minimizing that RMS. If its value were strictly less than $a$,
then the preceding chain would yield $|z_w(x_*)|<1/2$.
Hence the minimum RMS equals its quadratic mean $a$, so
\[
 \frac{\|W^1\bar x\|_2}{\sqrt N}=a
 \quad\text{for every Boolean }x.
\]
For every input, the entire output bound is consequently an equality:
$|z_w(x)|=F=1/2$. In particular the first hidden vector has the fixed nonzero
norm $\sqrt N\tanh a$.
If $L=2$, equality in readout Cauchy--Schwarz says that all first hidden
vectors are collinear with the same nonzero readout vector.
If $L\ge3$, equality in
$\|W^2h^1(x)\|_2\le\|W^2\|_F\|h^1(x)\|_2$ says the same: every nonzero row
of $W^2$ is collinear with $h^1(x)$, and $W^2\ne0$.
Thus, fixing any reference input $x_0$, there is
$\epsilon(x)\in\{-1,1\}$ with
\[
 h^1(x)=\epsilon(x)h^1(x_0).
\]
Oddness and injectivity of tanh imply
$W^1\bar x=\epsilon(x)W^1\bar x_0$.
Choose a nonzero coordinate of $W^1\bar x_0$; the corresponding first-layer
affine form therefore has constant nonzero absolute value on the cube.
An affine form $\alpha_0+\sum_i\alpha_i x_i$ with constant squared value can
have only one nonzero coefficient: expanding its square forces
$\alpha_0\alpha_i=0$ and $\alpha_i\alpha_j=0$ for $i\ne j$.
Therefore $\epsilon$ is a signed constant or dictator.
Every remaining layer is odd, so $z_w(x)=\epsilon(x)z_w(x_0)$, proving that
its Boolean function belongs to $\mathcal M_d$.

The strict-margin infimum causes no missing equality case.
For each representable Boolean function $g$, replace strict constraints by
$g(x)z_w(x)\ge1/2$. The relaxed feasible set is closed and its minimum norm
is attained by coercivity in the finite-dimensional parameter space.
Its minimum equals the strict infimum: scaling the readout by $1+\eta$ gives
strict feasibility with norm converging to the relaxed value as $\eta\downarrow0$.
Consequently every representable $g\notin\mathcal M_d$ has
$\kappa_N(g)>K_*$.

Finally, the small-variance Gaussian calculation in
Section~\ref{sec:strong-limit}, applied to the open strict-margin
implementation sets $U_{N,g}$, gives
\[
 \frac{\mu_{N,\sigma}(U_{N,g})}
      {\mu_{N,\sigma}(U_{N,f})}\longrightarrow0
 \quad\text{when }\kappa_N(g)>\kappa_N(f).
\]
There are finitely many Boolean functions. Therefore, for $f\in\mathcal M_d$,
\[
 \frac{\pi_{N,\sigma}(\mathcal M_d^c)}{\pi_{N,\sigma}(f)}\longrightarrow0.
\]
Choose $\sigma_0$ so that this ratio is at most $\varepsilon/2$.
At the displayed sample size, a union bound shows that, with probability
$1-\delta$, every consistent member of $\mathcal M_d$ has error at most
$\varepsilon/2$. The posterior mass outside $\mathcal M_d$ is at most
$\pi_{N,\sigma}(\mathcal M_d^c)/\pi_{N,\sigma}(f)\le\varepsilon/2$,
because $f$ is always consistent. Averaging errors proves the claimed
posterior risk guarantee.
\end{proof}

\paragraph{Scope.}
This is a result for the actual independent Gaussian prior in the strong-prior
limit, with the all-input margin conditioning retained. It does not settle
arbitrary circuit-simple targets, the specified nonzero variances in the
$\gamma\ge1$ papers, or a width/depth-uniform variance threshold.
It shows specifically why a uniform-truth-table low-norm prior counterexample
with a constant target cannot simply be transferred to this Gaussian prior.

% Optional corollary to gaussian-simple-targets.tex.
\begin{corollary}
For every fixed $d,N,L$, the globally margin-conditioned Gaussian function
prior converges, as $\sigma\downarrow0$, to the uniform distribution on
$\mathcal M_d=\{\pm1,\pm x_1,\ldots,\pm x_d\}$.
For data labeled by any $f\in\mathcal M_d$, the posterior converges to the
uniform distribution on the members of $\mathcal M_d$ consistent with the
data. In fact this posterior convergence is uniform over all such labeled
samples, including arbitrary sample sizes.
For uniform IID inputs, if
\[
 n\ge\log_2\frac{2d+1}{\delta},
\]
then with probability at least $1-\delta$ the limiting posterior is the point
mass at $f$.
\end{corollary}

\begin{proof}
The norm theorem shows that the prior mass outside $\mathcal M_d$ tends to
zero relative to the mass of any one member. Their masses are exactly equal
at every positive variance. Indeed, swapping columns $0$ and $i$ in the first
weight matrix preserves its independent Gaussian law. Let $T_i$ be the
involution of the Boolean cube defined by
\[
 (T_i x)_i=x_i,\qquad (T_i x)_j=x_i x_j\quad(j\ne i).
\]
Since the entire network, as a function of its augmented input, is odd, this
column swap transforms its output into
\[
 z_{w'}(x)=x_i z_w(T_i x).
\]
It therefore bijects the strict-margin implementation sets for $+1$ and
$+x_i$. Negating the readout preserves its Gaussian law and supplies the
negative functions. Thus all $2(d+1)$ implementation sets have equal mass.

Conditioning on any sample labeled by $f\in\mathcal M_d$ retains at least
$f$. If $\pi$ is the function prior, the posterior mass outside
$\mathcal M_d$ is bounded by $\pi(\mathcal M_d^c)/\pi(f)$, independently of
the sample and its size. Within $\mathcal M_d$, consistency conditioning is
exactly uniform. This proves the posterior convergence claim.
Under uniform inputs, each distinct $g\in\mathcal M_d$ disagrees with $f$ on
at least half the cube. Hence its probability of surviving $n$ labels is at
most $2^{-n}$. A union bound over the at most $2d+1$ competitors proves the
last assertion.
\end{proof}

% Optional proposition following the scalar characterization in gaussian-simple-targets.tex.
\begin{proposition}[The common depth cost]
For $d>0$ and integer $L\ge2$, let $K_*(d,L)$ be the minimum of
$da^2+\sum_{s=2}^{L-1}b_s^2+c^2$ over nonnegative scalar weights
subject to
$r_1=\tanh a$, $r_s=\tanh(b_sr_{s-1})$, and $cr_{L-1}\ge1/2$.
This is the energy of a constant-channel margin-$1/2$ realization.
Put $D=\max\{d,4\}$. Then
\[
 L(d/4)^{1/L}\ \le\ K_*(d,L)\ \le\ (L+12)(D/4)^{1/L}.
\]
In particular, for every fixed input dimension, $K_*(d,L)=L+O_d(1)$ as
$L\to\infty$. More precisely,
\[
 \log(d/4)\le\liminf_{L\to\infty}(K_*(d,L)-L)
 \le\limsup_{L\to\infty}(K_*(d,L)-L)
 \le12+\log(D/4).
\]
\end{proposition}

\begin{proof}
Since $\tanh t\le t$ for $t\ge0$, scalar feasibility implies
$ca\prod_{s=2}^{L-1}b_s\ge1/2$. The product of the $L$ weighted cost terms
$da^2,b_2^2,\ldots,b_{L-1}^2,c^2$ is at least $d/4$.
The arithmetic--geometric mean inequality proves the lower bound.

For the upper bound set
\[
 q=(D/4)^{1/(2L)}\ge1,\qquad
 k_s=\min\{s,L-s\},\qquad
 r_s=\frac{q^s}{\sqrt D\,k_s}\quad(1\le s<L).
\]
Then $0<r_s\le1/(2qk_s)\le1/2$ and $r_{L-1}=1/(2q)$.
Take
\[
 a=\operatorname{atanh}(r_1),\qquad
 b_s=\frac{\operatorname{atanh}(r_s)}{r_{s-1}},\qquad c=q.
\]
This realizes the hidden amplitudes $r_s$ exactly and has output $1/2$.
For $0<t\le1/2$, the power series gives
\[
 \frac{\operatorname{atanh}t}{t}
 \le1+\frac{t^2}{3(1-t^2)}\le1+\frac49t^2,
 \qquad
 \left(\frac{\operatorname{atanh}t}{t}\right)^2\le1+t^2.
\]
(The last step uses $t^2\le1/4$.) Hence
\[
 da^2\le dr_1^2+dr_1^4\le q^2+\frac{q^4}{D}\le q^2+\frac14,
 \qquad c^2=q^2.
\]
Put $R_s=k_{s-1}/k_s$ for $2\le s<L$.
The ratios occur in reciprocal pairs $k/(k+1),(k+1)/k$, with at most one
additional ratio equal to one. Each pair contributes at most $4/k^2$ above
two, since
\[
 \left(\frac{k+1}k\right)^2+
 \left(\frac{k}{k+1}\right)^2-2
 =\left(\frac{k+1}k-\frac{k}{k+1}\right)^2\le\frac4{k^2}.
\]
Therefore $\sum_{s=2}^{L-1}R_s^2\le L-2+4\sum_{k\ge1}k^{-2}<L-2+7$.
Moreover $R_s\le2$ and $q^{2s+2}\le q^{2L}=D/4$, so
\[
 \begin{split}
 \sum_{s=2}^{L-1}b_s^2
 &\le q^2\sum_{s=2}^{L-1}R_s^2
       +\sum_{s=2}^{L-1}q^2R_s^2r_s^2\\
 &\le q^2(L-2+7)+\sum_{s=2}^{L-1}\frac1{k_s^2}
 <q^2(L-2+7)+4.
 \end{split}
\]
Adding the endpoint costs gives
$K_*(d,L)\le q^2(L+7)+17/4\le q^2(L+12)$.
The argument also applies at $L=2$, when the internal sums are empty.
Finally, expand $L e^{\log(d/4)/L}$ and
$(L+12)e^{\log(D/4)/L}$ for fixed $d$ to obtain the asymptotic statements.
\end{proof}

% Standalone appendix fragment; amsmath, amssymb, amsthm required.
\section{What excess Gaussian energy does and does not control}
\label{app:excessenergy}

Write a width-$N$, depth-$L$ tanh network as
$h^1=\tanh(W^1\bar x)$, $h^s=\tanh(W^sh^{s-1})$ for $2\le s<L$,
and $z=v^Th^{L-1}$, where $\bar x=(1,x)$ and $x\in\{-1,1\}^d$.
Its weighted Gaussian energy is
\[
 \kappa_N=\frac dN\|W^1\|_F^2+
       \sum_{s=2}^{L-1}\|W^s\|_F^2+N\|v\|_2^2.
\]
The following facts concern this actual energy. They do not establish a
sample-complexity lower bound for the unrestricted Gaussian posterior.

\begin{proposition}[Excess energy controls total deviation from rank one]
Suppose $|z(x)|\ge1/2$ on the whole Boolean cube. Let $R_s$ be a best
rank-one Frobenius approximation to $W^s$, and let
$\rho_s=\|W^s\|_F^2/\|W^s\|_{\mathrm{op}}^2$ be its stable rank.
Then
\begin{align}
 \sum_{s=2}^{L-1}\|W^s-R_s\|_F^2
 &\le \kappa_N-L(d/4)^{1/L},\label{eq:excessrank}\\
 \sum_{s=2}^{L-1}\log\rho_s
 &\le L\log(\kappa_N/L)-\log(d/4).
\end{align}
In particular, if $\kappa_N\le L+B$, both right-hand sides are at most
$B-\log(d/4)$.
\end{proposition}

\begin{proof}
Set $a=\|W^1\|_F/\sqrt N$, $c=\sqrt N\|v\|_2$, and
$p_s=\|W^s\|_{\mathrm{op}}$. Cube averaging supplies an input with
$\|W^1\bar x\|_2/\sqrt N\le a$. Tanh contraction gives
$1/2\le ac\prod_s p_s$. Thus the product of the $L$ numbers
$da^2,p_2^2,\ldots,p_{L-1}^2,c^2$ is at least $d/4$, and their sum is at
least $L(d/4)^{1/L}$. Subtract this sum from $\kappa_N$ and use
$\|W^s-R_s\|_F^2=\|W^s\|_F^2-p_s^2$ to obtain
\eqref{eq:excessrank}. Alternatively, the product of the Frobenius energy
terms is at least $(d/4)\prod_s\rho_s$ and at most $(\kappa_N/L)^L$,
which proves the logarithmic bound. Finally use $e^u\ge1+u$ and
$\log(1+u)\le u$.
\end{proof}

Thus bounded excess permits only a bounded total squared residual after
replacing internal matrices by rank-one matrices. This does not justify
replacing those matrices in the computation: the suffix can amplify the
replacement errors. The next example shows that even a constant excess
allowance does not imply inverse-polynomial hidden-state robustness.

\begin{proposition}[Bounded excess permits a fragile scalar transmission]
For $d=4$, every even $L\ge4$, and every $N\ge1$, there is a network with
$z(x)=1$ for all inputs and
\[
 \kappa_N\le L+76\le K_*(4,L)+76,
\]
whose output can be made zero by a perturbation at one hidden layer of
normalized RMS at most $L^{-1}e^{-\sqrt L/4}$.
Here $K_*(4,L)$ is the scalar minimum for margin-$1/2$ transmission
defined above. It also equals the globally minimum margin-$1/2$ energy
for this model, by the simple-target result in the research report.
\end{proposition}

\begin{proof}
Put $n=L/2$, $h=L^{-1/2}$, $k_s=\min(s,L-s)$, and
\[
 r_s=\frac{e^{-h(k_s-1)}}{2k_s},\qquad 1\le s<L.
\]
Give each first-layer row the sole constant-input weight
$\operatorname{atanh}(1/2)$. Let every entry of $W^s$ be $b_s/N$, where
$b_s=\operatorname{atanh}(r_s)/r_{s-1}$, and let $v_i=2/N$.
The hidden vector is exactly $r_s\mathbf1$, so the output is
$2r_{L-1}=1$ and the energy is
\[
 4\operatorname{atanh}(1/2)^2+\sum_{s=2}^{L-1}b_s^2+4.
\]
For completeness, the following loose constants are uniform in depth.
The ratios $r_s/r_{s-1}$ occur in reciprocal pairs $t_k,t_k^{-1}$,
where $t_k=(1+1/k)e^h$ and $1\le k<n$. Since $h\le1/2$,
\[
 t_k-t_k^{-1}\le e^h(2/k+2h),\qquad
 t_k^2+t_k^{-2}-2\le8e^{2h}(k^{-2}+h^2).
\]
Using $e^{2h}<3$, $\sum_{k\ge1}k^{-2}<2$, and $(n-1)h^2\le1/2$ gives
$\sum_s(r_s/r_{s-1})^2<L-2+60$.
Also $\max_s(r_s/r_{s-1})^2<12$ and $\sum_s r_s^2<1$.
For $0<t\le1/2$ the power series gives
$(\operatorname{atanh}(t)/t)^2\le1+t^2$. Consequently
\[
 \sum_s b_s^2<L-2+72,
 \qquad 4\operatorname{atanh}(1/2)^2\le5/4,
 \qquad \kappa_N<L+76.
\]
The baseline lower bound gives $K_*(4,L)\ge L$.
Finally, adding $-r_n\mathbf1$ after the middle activation zeros that
state and hence every later state. Its normalized RMS is
\[
 r_n=L^{-1}\exp\!\left(-\frac{L/2-1}{\sqrt L}\right)
 \le L^{-1}e^{-\sqrt L/4}.
\]
\end{proof}

This example represents a statistically easy target: the report's
simple-target theorem gives depth-independent strong-prior learning
for constants, at fixed width and depth before taking the variance limit. It therefore
separates fragility of a cheap parameter realization from difficulty of
learning its Boolean function. It does not show that the fragile realization
is a minimum-energy representative.

\begin{proposition}[Additional depth has asymptotically unit energy cost]
Suppose a width-$N$, depth-$L_0$ network realizes $f$ with margin at least
$1/2$. Write $c^2=N\|v\|_2^2>0$ and $E_{\mathrm{pre}}=\kappa_N-c^2$.
For every integer $t\ge1$, there is a width-$N$, depth-$(L_0+t)$ realization
of the same function, with margin at least $1/2$ and energy at most
\[
 E_{\mathrm{pre}}+(t+13)\max\{c^2,1\}^{1/(t+1)}.
\]
Write $\kappa_{N,L}(f)$ for the infimum energy of depth-$L$, width-$N$
representations with signed output strictly greater than $1/2$ on the cube.
For each fixed representable $f,d,N$, this infimum satisfies
$\kappa_{N,L}(f)=L+O_{f,d,N}(1)$ as $L\to\infty$.
\end{proposition}

\begin{proof}
Extend the scalar minimum $K_*(D,m)$ to any real $D>0$: minimize
$Da^2+\sum_{j=2}^{m-1}b_j^2+c'^2$ subject to $c'r_{m-1}\ge1/2$,
where $r_1=\tanh a$ and $r_j=\tanh(b_jr_{j-1})$.
The scalar baseline upper-bound argument works unchanged for real $D>0$:
\[
 K_*(D,m)\le(m+12)(\max\{D,4\}/4)^{1/m}.
\]
Indeed, setting $q=(\max\{D,4\}/4)^{1/(2m)}$ and
$r_j=q^j/[\sqrt{\max\{D,4\}}\min(j,m-j)]$ gives this bound by
$a=\operatorname{atanh}r_1$, $b_j=\operatorname{atanh}(r_j)/r_{j-1}$,
$c'=q$, and the summable reciprocal-ratio estimates used earlier.

Take the scalar tuple for $D=4c^2$ and $m=t+1$. Replace the old readout
by a hidden layer whose every row is $2a v^T$, followed by the remaining
scalar layers and readout. Its first new preactivation is $2az(x)$, whose
absolute value is at least $a$ and whose sign is $f(x)$. Monotonicity and
oddness give the same sign with final margin at least $1/2$.
The first new block has energy $4a^2c^2$; the remaining new blocks have
exactly the remaining scalar costs. The total is therefore
$E_{\mathrm{pre}}+K_*(4c^2,t+1)$, proving the displayed upper bound.
An arbitrarily small readout increase supplies strict margins.
Expanding this bound as $t\to\infty$, and combining it with
$L(d/4)^{1/L}\le\kappa_{N,L}(f)$, proves the final assertion.
\end{proof}

These energy estimates alone do not settle unrestricted Gaussian sample
complexity. They rule out exponential growth of minimum target energy under depth
padding, while showing that bounded excess does not make every cheap
realization robust. A target-relative argument must also control competing
functions and their prior masses; subtracting a common transmission cost
alone does not do so.

\section{Robust posterior compression uniformly for \texorpdfstring{$\gamma\ge1$}{gamma at least one}}
\label{sec:robust-entropy}

This section compresses a distribution of robust networks into a fixed
finite class. It uses robustness of the original network while constructing
a candidate, and tracks entropy
only in the original weights that have not yet been processed. Consequently
the intermediate compressed prefix need not have a density under a Gaussian
prior. The result applies to a posterior containing many different functions.

Let $\mu_{N,t}$, $0<t\le1$, have independent weight-entry variances
$t^2/d$, $t^2/N$, and $t^2/N^2$ in the first, internal, and readout layers.
The exponent-$\gamma$ prior is obtained by $t=N^{(1-\gamma)/2}$.
All-layer sign robustness refers to the total noise norm
in~\eqref{eq:noisenorm}; the clean $1/2$ margin may additionally be imposed.

\begin{theorem}[Uniform robust posterior compression]
\label{thm:robust-entropy}
Fix a scalar $1$-Lipschitz activation $|\phi|\le1$, and parameters
$d\ge1$, $L\ge2$, $K\ge1$, $0<r\le1$, and $0<\alpha\le1/8$.
There is a finite class $\mathcal H=\mathcal H(d,L,K,r,\alpha)$, consisting of
signs of bounded-weight depth-$L$ networks, such that
\[
 \log|\mathcal H|\le\poly(d,L,K,r^{-1},\alpha^{-1}).
\]
The polynomial is universal and does not depend on $N$ or $t$.
If $Q\ll\mu_{N,t}$ is supported on all-layer sign-$r$ robust networks and
\[
 D_{\rm KL}(Q\Vert\mu_{N,t})\le NK/t^2,
\]
then there is a distribution $\Gamma_Q$ on $\mathcal H$ such that, for every
$x\in\X$, a coupling of $w\sim Q$ and $h\sim\Gamma_Q$ satisfies
\[
 \Pp\{h(x)\ne\sgn z_w(x)\}\le\alpha.
\]
The coupling may depend on $x$; its two marginals do not. Neither the support
of $Q$ nor its samples need represent a single Boolean function.
\end{theorem}

We give the proof, including deliberately loose polynomial choices.
Two elementary entropy facts organize the construction. First, for a standard
$D$-dimensional Gaussian $G$, the entropy variational inequality gives
\begin{equation}\label{eq:entropy-moment}
 \E_P\|Z\|^2\le4D_{\rm KL}(P\Vert G)+2D\log2.
\end{equation}
Second, Gaussian Talagrand~\cite{talagrand} gives
$W_2^2(P,G)\le2D_{\rm KL}(P\Vert G)$.

\begin{lemma}[Entropy of the unprocessed weights]
\label{lem:future-entropy}
Let $A,X,Y$ be random variables, with fixed reference laws $P_X,P_Y$, and put
\[
 B=\E_A D_{\rm KL}(Q_{X,Y\mid A}\Vert P_X\otimes P_Y).
\]
Let $J$ be a finite message, possibly depending on $(A,X,Y)$, and let
$H=f(A,J,X)$. If $A'$ is a randomized function of $(A,J,H)$ whose fresh
randomness is conditionally independent of $Y$, then
\[
 \E_{A'}D_{\rm KL}(Q_{Y\mid A'}\Vert P_Y)\le B+H(J\mid A).
\]
\end{lemma}
\begin{proof}
The left side equals $D_{\rm KL}(Q_Y\Vert P_Y)+I(Y;A')$.
Data processing and
\[
 I(Y;J,H\mid A)\le I(Y;X,J\mid A)
 \le I(Y;X\mid A)+H(J\mid A)
\]
bound it by $\E_{A,X}D_{\rm KL}(Q_{Y\mid A,X}\Vert P_Y)+H(J\mid A)$,
which is at most the claimed bound by the KL chain rule. These expressions
can equivalently be justified by finite partitions and lower semicontinuity.
\end{proof}

Here $A$ will be a constructed prefix, $X$ the current original weight block,
and $Y$ all later original blocks. The lemma imposes no reference law on $A$.

\begin{lemma}[Whole-layer active directions]
\label{lem:block-entropy}
Condition on $(A,Y)$ and let $X$ be a standardized Gaussian-reference matrix,
with conditional relative entropy $R(A,Y)$. For vectors $b_x(A)$ of norm at
most one, one can select a list $J$ of inputs and its span
$U=\operatorname{span}\{b_x:x\in J\}$ with
\[
 \dim U\le\frac{t^2R(A,Y)}{N\delta}
\]
such that preserving $H=tXU$ and replacing the orthogonal coordinates by fresh
Gaussian coordinates changes the physical evaluation vectors $tXb_x$, under
a suitable coupling, by expected squared RMS at most $2\delta$ at each $x$.
The candidate matrix law preserves $(A,Y,H)$ and does not depend on $x$.
\end{lemma}
\begin{proof}
Use a fixed ordering of the finite input set. Whenever the averaged conditional
KL of the entire residual evaluation vector $Xv$ exceeds $N\delta/t^2$, add
the unit residual direction $v$ of that input. Each query is selected from
the conditional law given $(A,Y)$, averaging over active-coordinate values;
thus $J$ is a function of $(A,Y)$. Orthogonal Gaussian coordinates
are independent, so the active-coordinate entropy increases by precisely this
conditional KL. The displayed rank bound follows. At termination, Gaussian
Talagrand on the entire evaluation vector, multiplication by the squared
residual-vector length, and division by $N/t^2$ give the error bound.
Scalar-vector couplings extend to full matrices by disintegration.
The basis of $U$ is a measurable function of $A,J$ under fixed Gram--Schmidt;
it has no further dependence on $Y$ once those variables are supplied.
\end{proof}

\paragraph{Parameter choices.}
Write
\begin{align*}
 e&=\alpha r^2/(1000L),& \eta&=\alpha/(100L),&
 q&=\lceil2K/(e\eta)\rceil,\\
 D_0&=[(d+1)q+2]\log2,& V&=8K+2q,& B_H^2&=4V/e,\\
 T&=128(q+1)(B_H^2+1)/e,& a&=e/[128(B_H^2+1)],\\
 K_0&=(4K+2\log2)/\eta,& B_v&=16K_0/r.
\end{align*}
Choose $m$ to be the least positive integer no smaller than both
\begin{align*}
 &10000T^4(q+1)^2a^{-2}[d+\log(10(q+1)/\eta)+1],\\
 &10000B_v^2r^{-2}[d+\log(10/\eta)+1],
\end{align*}
and put
\[
 N_0=\max\{m,\lceil LD_0/K\rceil,1\},\quad
 B_G=\max\{2,2\sqrt{\log(3m/e)}\},\quad
 B_1^2=4[4K+2(d+1)\log2]/e.
\]
All quantities have a fixed polynomial upper bound in the theorem parameters.

\paragraph{The small-width case.}
If $N\le N_0$, let
$P_0=N_0(d+1)+(L-2)N_0^2+N_0$. Applying~\eqref{eq:entropy-moment} to all
standardized weights, and then multiplying by $t^2$, gives
\[
 \E_Q\|tZ_{\rm all}\|^2\le4N_0K+2P_0\log2.
\]
With probability at least $1-\alpha/4$, all actual weights are bounded by
\[
 B_{\rm small}=[4(4N_0K+2P_0\log2)/\alpha]^{1/2}.
\]
Keep the original network on this event and use a default classifier otherwise.
Padding with nodes having zero outgoing weights preserves the original output.
The successful candidate has output margin greater than $r$.
It remains to treat $N>N_0$.

\paragraph{The virtual prefix and its entropy budget.}
The constructed prefix $A$ has $N$ virtual nodes in its current layer. Each
virtual row acts on $m$ stored nodes of the preceding layer. The auxiliary
description may keep all $N$ rows; the final network will keep just the rows
selected at the following stage. Track the conditional KL of the unprocessed
\emph{original} weights. By Lemma~\ref{lem:future-entropy}, it is at most
\[
 B_s\le NK/t^2+sD_0\le2NK/t^2.                              \tag{*}
\]
This is not a KL bound for the singular compressed prefix itself.

For the first layer, zero every physical standardized row $tZ_{1,i}$ of norm
greater than $B_1$, and retain all other rows. Bounded activations and
\eqref{eq:entropy-moment} bound the expected squared RMS activation change
at each input by $e$. This deterministic prefix channel does not increase the
future entropy budget.

\paragraph{An internal layer: retain active data.}
Condition on the current prefix $A$ and the future original weights $Y$.
Apply Lemma~\ref{lem:block-entropy} to
$b_x=h_A(x)/\sqrt N$ and $\delta=e$. Abort if more than $q$ inputs are needed.
By (*) and Markov, this failure has probability at most $\eta$.
The successful query list, together with an abort flag, has a code of at most
$[(d+1)q+2]$ bits. Its entropy is at most $D_0$.

Let $H=tXU$ be the physical active matrix. Averaging the Gaussian entropy
moment bound, with the success indicator rather than conditioning on success, gives
\[
 \E\left[\frac{\|H\|_F^2}{N};\ \text{success}\right]
 \le4t^2B_s/N+2t^2q\log2\le V.
\]
Zero every row $H_i$ whose norm exceeds $B_H$. The expected squared RMS
activation change from this operation is at most $4V/B_H^2=e$.
All rows are clipped before any later selection, including selections that
may depend on future weights.

\paragraph{An internal layer: reconstruct using $m$ source nodes.}
For the orthonormal active basis let $u(j)=\sqrt N\,U_{j,:}$; uniform source
sampling gives $\E_j u(j)u(j)^\top=I$.
Clip the \emph{whole vector} at radius $T$:
$c(j)=u(j)\mathbf1_{\{\|u(j)\|\le T\}}$. Then
\[
 \E cc^\top\preceq I,\qquad \E\|u-c\|\le q/T.
\]
Sample $m$ uniform source indices $S_j$. Entrywise Hoeffding and a union bound
give, with probability at least $1-\eta$,
\begin{align*}
 \|\widehat G-\E cc^\top\|_{\op}&\le a,\quad
 \sup_x\|\widehat v(x)-\E c h_A(x)\|\le a,\\
 \sup_{x,y}|\widehat K(x,y)-K(x,y)|&\le a,
\end{align*}
where $\widehat G=m^{-1}\sum_jc(S_j)c(S_j)^\top$,
$\widehat v(x)=m^{-1}\sum_jc(S_j)h_{A,S_j}(x)$, and $\widehat K$ is the
empirical feature covariance; its population counterpart is
$K(x,y)=N^{-1}\sum_i h_{A,i}(x)h_{A,i}(y)$.
Abort on failure; this event depends only on
$A,J,S$, so it does not reveal further future information.

Set $\widetilde c_j=c(S_j)/\sqrt{1+a}$, and let $D$ have rows
$\widetilde c_j^\top/\sqrt m$. Then $D^\top D\preceq I$ and
$I-DD^\top\succeq0$. Independently for every target virtual row draw
$\xi_i\sim N(0,I-DD^\top)$ and use coefficients
\begin{equation}\label{eq:finite-row}
 \widehat w_{ij}=m^{-1}H_i^{\rm clip}\cdot\widetilde c_j
                         +t\xi_{ij}/\sqrt m.
\end{equation}
These are actual neural-network rows on the $m$ sampled nodes.
Their active feature vector is
$\widehat\psi(x)=m^{-1}\sum_j\widetilde c_jh_{A,S_j}(x)$, and the lazy
covariance is exactly $t^2(\widehat K-\widehat\psi^\top\widehat\psi)$.
For the ideal field, write $\psi(x)=U^\top b_x$ and
$C(x,y)=K(x,y)-\psi(x)^\top\psi(y)$. Direct calculation gives
\[
 \sup_x\|\widehat\psi(x)-\psi(x)\|\le q/T+2a=:\zeta,
 \qquad \sup_{x,y}|\widehat C(x,y)-C(x,y)|\le a+2\zeta.
\]
Here $\widehat C=\widehat K-\widehat\psi^\top\widehat\psi$.
Bessel's inequality and $D^\top D\preceq I$ give
$\|\psi(x)\|,\|\widehat\psi(x)\|\le1$, which justify the covariance bound.
At a fixed input, coupling scalar Gaussian evaluations with a common standard
normal bounds the squared preactivation transport error, per row, by
\[
 B_H^2\zeta^2+t^2(a+2\zeta)\le e/8.
\]
Extend these scalar couplings to full row fields. The candidate weight law is
the same at every input; there is no claim of a single Gaussian-field coupling
that is close at all inputs simultaneously.

Finally clip \emph{all} coordinates $\xi_{ij}$ to $[-B_G,B_G]$ before later
row selection. Since each has variance at most one,
\[
 \E\left|\frac{t}{\sqrt m}\sum_j
   (\xi_{ij}-\operatorname{clip}_{B_G}\xi_{ij})h_{A,S_j}(x)\right|^2
 \le3m e^{-B_G^2/4}\le e.
\]
Thus every virtual coefficient is bounded by $B_HT/m+B_G/\sqrt m$,
without a union bound over $N$.

All newly constructed data are a randomized function of $A,J,H$.
Lemma~\ref{lem:future-entropy} therefore proves (*) for the next step.
Combining entropy replacement, active clipping, finite reconstruction, and
Gaussian clipping by $(a_1+\cdots+a_4)^2\le4\sum_i a_i^2$, the expected
squared RMS activation injection at this layer is less than $20e$.

\paragraph{The readout and final network.}
The original readout marginal has never changed. Its physical moment satisfies
$\E_Q[N\|v\|^2]\le4K+2\log2$.
Abort if $N\|v\|^2>K_0$, an event of probability at most $\eta$.
Zero coefficients with $|Nv_i|>B_v$. This changes the output uniformly on
bounded virtual feature vectors by at most $K_0/B_v=r/16$.
Sampling $m$ uniform last-layer sources and using readout coefficients
$(Nv_i\mathbf1_{|Nv_i|\le B_v})/m$ adds at most $r/16$ uniformly on $\X$,
except with probability $\eta$, by Hoeffding and the chosen $m$.

At each stage, retain in the final network just the $m$ source rows selected
from the preceding virtual layer. Each such row already acts on the previously
stored $m$ nodes. The result is an ordinary width-$m$ network; repeated indices
are duplicate nodes. Its weights are bounded by
\[
 B_{\rm large}=\max\{B_1,B_HT/m+B_G/\sqrt m,B_v/m,1\}.
\]
The larger virtual construction is not part of the final predictor.

\paragraph{The original network receives all perturbations.}
Fix an input $x$ and glue the conditional couplings, retaining the original
weight history. The current and future original blocks are untouched until
their own stages. Define
\[
 \eta^s=h^s_{\rm virtual}(x)
           -\phi(W^s_{\rm original}h^{s-1}_{\rm virtual}(x)).
\]
The analogous readout difference is $\eta^L$. Later stages leave earlier
computational prefixes unchanged, so the final candidate output is exactly
the \emph{original} network's output under these postactivation injections.
No intermediate robustness assumption is made.

For precision, the canonical construction draws its new Gaussian coefficients
independently of the discarded current block given $(A,J,H)$. The comparison
coupling may correlate them with that discarded block. It preserves the same
joint marginal of the new prefix and the future original blocks, as well as
the original full-network marginal, but need not preserve the canonical joint
law with all discarded weights. The next entropy step uses the first marginal;
original-network robustness uses the second. There is no conflicting demand
for independence in the comparison coupling itself.

More explicitly, the induction invariant is that the comparison law has
full original-weight marginal $Q$ and the canonical marginal of
$(A,X,Y)$. The conditional transports supply a joint law of
$(A,X,Y,A')$ with this incoming marginal and the prescribed outgoing
$(A',Y)$ marginal. Extend it by the previous conditional law of the past
original weights given $(A,X,Y)$. The extension keeps the complete original
weight law unchanged, preserves all earlier computational prefixes, and
establishes the same invariant for the next layer. This is ordinary
disintegration on finite-dimensional standard Borel spaces.

Abort probability is at most $2L\eta$. The unconditional moment bound gives
\[
 \Pp\left\{\text{no abort and }\sum_{s<L}\|\eta^s\|^2/N>r^2/8\right\}
 \le160Le/r^2=0.16\alpha.
\]
Otherwise the readout injection is at most $r/8$, and total squared noise is
at most $9r^2/64<r^2/4$. Robustness then gives candidate margin greater than
$r/2$ relative to the original network's label: a smaller margin could be
cancelled by an additional output shift within the radius-$r$ noise ball.
The exceptional probability is less than $\alpha/4$.

\paragraph{A fixed finite class.}
Pad candidates to width $M=N_0$ and take
$B=\max(B_{\rm small},B_{\rm large},1)$. Round every parameter to a fixed
grid on $[-B,B]$ of mesh at most
\[
 h=\frac{r}{16(d+L+1)(2M(B+1))^L}.
\]
The bounded-activation, Lipschitz layer recursion gives uniform output change
at most $r/8$. Thus rounding preserves every successful coupled prediction.
Writing $P=M(d+1)+(L-2)M^2+M$, the grid class satisfies
\[
 \log|\mathcal H|\le
 P\log\left[3+\frac{32B(d+L+1)(2M(B+1))^L}{r}\right].
\]
Although grid precision is exponentially small in $L$, its number of bits is
polynomial in $L$. All other bounds above are polynomial too. Rounding is a
fixed map on candidate networks, giving the common marginal $\Gamma_Q$ and
proving Theorem~\ref{thm:robust-entropy}.

\begin{corollary}[Polynomial-depth learning at every width]
\label{cor:robust-entropy-pac}
Suppose the raw Gaussian target event, consisting of networks with all-layer
sign margin $r$ and clean margin $1/2$ for $f$, has mass at least
$\exp(-NK/t^2)$. Restrict the prior to all sign-$r$ robust networks having
global clean margin $1/2$, regardless of their Boolean function, then condition
on $n$ noiseless iid training examples. With probability at least $1-\beta$,
the resulting posterior satisfies
\[
 \operatorname{GibbsErr}(Q)
 \le2\alpha+\sqrt{\frac{\log(2|\mathcal H|/\beta)}{2n}}.
\]
Consequently the sample bound is jointly polynomial in
$d,L,K,r^{-1}$ and inverse desired error, uniformly in $N$ and $0<t\le1$.
For posterior-confidence error at threshold $2/3$, multiply Gibbs error by
three; for threshold $1/2$, multiply it by two.
\end{corollary}
\begin{proof}
The posterior is the raw Gaussian conditioned on an event containing the target
event, so its KL is at most $NK/t^2$. Its training error is zero. The compression
couplings imply that $\Gamma_Q$ has expected empirical error at most $\alpha$.
Finite-class uniform convergence holds for every member of $\mathcal H$, hence
also for this sample-dependent mixture. Its true error is at most $\alpha$
plus the displayed square-root term. The pointwise coupling adds at most one
more $\alpha$ to recover original posterior Gibbs error. This argument does
not require a global coupling over inputs or a common function for all draws.
The confidence-error statements follow from Markov on posterior error
probability at a random test input.
\end{proof}

\section{A finite-radius certificate using clean Jacobian products}
\label{sec:clean-transport-certificate}

Products of individual spectral norms discard cancellations between layers.
The following sufficient test retains the exact clean transition products and
bounds the nonlinear remainder separately. It is a finite-radius statement,
not an inference from the infinitesimal Jacobian alone.

Write the $H=L-1$ hidden layers in column-vector convention as
$h_s=\tanh(W_sh_{s-1})$, with output $z=v^\top h_H$ and width $N$.
Fix an input and its clean preactivations $a_s=W_sh_{s-1}$. Define
\[
 A_s=\operatorname{diag}(\operatorname{sech}^2 a_s)W_s,
 \qquad P_{s,t}=A_s\cdots A_{t+1}\ (t<s),\quad P_{s,s}=I.
\]
The Gramian for normalized additive hidden disturbances is
\[
 G_1=NI,\qquad G_s=A_sG_{s-1}A_s^\top+NI.
\]
For a proposed total-noise radius $r$, recursively put $u_1=0$ and, for $s\ge2$,
\begin{align}
 b_s&=r\sqrt{\operatorname{diag}(W_sG_{s-1}W_s^\top)}
       +\sum_{t=2}^{s-1}|W_sP_{s-1,t}|u_t,\label{eq:transport-b}\\
 u_s&=\frac{2}{3\sqrt3}\,b_s^{\odot2}.\label{eq:transport-u}
\end{align}
Absolute values, the square root, and the square in these expressions are
coordinatewise. Finally define
\begin{equation}\label{eq:transport-output}
 B(r,x)=r\sqrt{v^\top G_Hv+1}
                  +\sum_{t=2}^H|v^\top P_{H,t}|u_t.
\end{equation}

\begin{proposition}[Finite-radius transport certificate]
Every simultaneous postactivation and readout disturbance satisfying
$\sum_{s=1}^H\|\eta_s\|^2/N+\eta_{\rm out}^2\le r^2$ changes the output
at $x$ by at most $B(r,x)$. Consequently, if $B(r,x)<|z(x)|$ on the full
input cube, the network preserves its own clean signs throughout this ball.
\end{proposition}
\begin{proof}
Let $\Delta h_s$ be the hidden displacement and write its exact recursion as
\[
 \Delta h_s=A_s\Delta h_{s-1}+\eta_s+R_s,
 \qquad
 \Delta h_s=\sum_{t=1}^sP_{s,t}\eta_t+\sum_{t=2}^sP_{s,t}R_t.
\]
There is no input perturbation, so $R_1=0$. The global scalar bound
$|\tanh''|\le4/(3\sqrt3)$ implies
$|R_s|\le[2/(3\sqrt3)]|W_s\Delta h_{s-1}|^{\odot2}$.
Assuming $|R_t|\le u_t$ at preceding layers, Cauchy--Schwarz applied to the
linear disturbance map, followed by the coordinatewise triangle inequality
for the remainder terms, gives $|W_s\Delta h_{s-1}|\le b_s$.
This proves $|R_s|\le u_s$ inductively. At the readout, Cauchy--Schwarz
including scalar output noise gives the first term of~\eqref{eq:transport-output};
the remainder sum gives the second. No statistical independence or
small-noise approximation is used.
\end{proof}

The test can be sharpened by replacing the global curvature bound in
\eqref{eq:transport-u} with the supremum of $|\tanh''|$ on
$[a_{s,j}-b_{s,j},a_{s,j}+b_{s,j}]$. An alternative exact scalar remainder
envelope evaluates
$\tanh(a+q)-\tanh(a)-\operatorname{sech}^2(a)q$
at $q=\pm b$ and at $q=-2a$ when $|2a|\le b$: these and $q=0$ exhaust
its extrema on the interval. The global-curvature version above already
suffices for the numerical improvement below.

\paragraph{Application to the saved models.}
We evaluated the exact support condition
$\min_x|z(x)|>1/2$ and sign preservation at
$r=.01/\sqrt{1+S_N(w)}$ on all 36 saved width-32 models, using all 64
inputs of the six-bit cube. The preceding spectral/derivative-range test established
membership for 28 models, excluded three by their clean absolute margin,
and left five unresolved. The new test settles all five:
\begin{center}
\begin{tabular}{lrrrr}
\toprule
Task & Hidden layers, seed & Required $r$ & $\max_x B(r,x)$ & $\min_x|z(x)|$\\
\midrule
Parity & $16,2$ & $.000882$ & $.012663$ & $.996070$\\
Majority & $64,1$ & $.000943$ & $.004762$ & $.999996$\\
Majority & $64,2$ & $.000908$ & $.006023$ & $.999993$\\
Parity & $32,2$ & $.001018$ & $.010793$ & $.998686$\\
Parity & $64,0$ & $.000872$ & $.007722$ & $.999419$\\
\bottomrule
\end{tabular}
\end{center}
Thus all 33 successful fits pass a sufficient test for this precise support
condition. The three remaining models are excluded because their minimum
absolute clean outputs are $.28125$, $.03125$, and $.06250$, respectively.
Their target misclassification alone was not used to infer failure of
own-label robustness. The original cohort has $24/24$ models passing; the
separate depth-32/64 stress cohort has $9/12$ passing and three excluded.
No run was omitted, retrained, or subjected to new attacks for this check.

\paragraph{Numerical and empirical qualification.}
Two independently inspected implementations agree: one uses the Gramian and
local curvature, while the other concatenates the exact clean disturbance
transfer matrices and uses the global curvature bound. Clean linear gains
also agree with independent backpropagation to relative discrepancy below
$1.2\times10^{-15}$. All computations use float64, including spectral norms;
they are numerical evaluations of rigorous sufficient tests, not
interval-verified exact-real certificates. The earlier audit and all saved
weights are retained. These finite examples demonstrate applicability to
some gradient-trained networks through 64 hidden layers; they establish
neither asymptotic polynomial depth scaling nor a universal property of SGD.


\begin{thebibliography}{9}
\bibitem{conditional2026} ChatGPT-6 Astra, K. H\"anni, and D. Vaintrob.
\emph{A conditional obstruction to polynomial sample complexity in unrestricted NNbayes.}
29 October 2026.
\href{https://github.com/kaarelh/AI-works/tree/main/neural-network-generalization/works-in-progress/conditional-nonpolynomial-learning}{Source and PDF}.
\bibitem{aboveone} ChatGPT-6 Astra, Claude Fable 5 and Opus 5,
K. H\"anni, and D. Vaintrob.
\emph{Width-uniform Bayesian generalization above the mean-field scaling.}
2026.
\url{https://github.com/kaarelh/AI-works/tree/main/neural-network-generalization/main-results/gamma-above-one}.
\bibitem{fable} Claude Fable 5. \emph{Robustness kills the exp(depth):
subsampling noise-stable nets at width $\poly(d,K,1/\tau)(d+\log L)$.}
27 August 2026. Unpublished research note, \texttt{robust-depth.pdf}.
\bibitem{project} ChatGPT-6 Astra, Claude (Fable 5, Opus 5, and Opus 5.5), K. H\"anni, and D. Vaintrob.
\emph{Bayesian neural-network generalization: main papers and Lean project.}
\url{https://github.com/kaarelh/AI-works/tree/main/neural-network-generalization}.
\bibitem{weima2020} C. Wei and T. Ma.
\emph{Improved Sample Complexities for Deep Networks and Robust Classification
via an All-Layer Margin.} ICLR 2020. \url{https://arxiv.org/abs/1910.04284}.
\bibitem{weima2019} C. Wei and T. Ma.
\emph{Data-dependent Sample Complexity of Deep Neural Networks via Lipschitz
Augmentation.} NeurIPS 2019. \url{https://arxiv.org/abs/1905.03684}.
\bibitem{arora} S. Arora, R. Ge, B. Neyshabur, and Y. Zhang.
\emph{Stronger Generalization Bounds for Deep Nets via a Compression Approach.}
ICML 2018. \url{https://proceedings.mlr.press/v80/arora18b.html}.
\bibitem{weichenma} C. Wei, Y. Chen, and T. Ma.
\emph{Statistically Meaningful Approximation: a Case Study on Approximating
Turing Machines with Transformers.} NeurIPS 2022.
\url{https://papers.nips.cc/paper_files/paper/2022/file/4ebf1d74f53ece08512a23309d58df89-Paper-Conference.pdf}.
\bibitem{tarres} P. Tarr\`es and Y. Yao.
\emph{Online Learning as Stochastic Approximation of Regularization Paths.}
Appendix A. \url{https://arxiv.org/abs/1103.5538}.
\bibitem{talagrand} M. Talagrand.
\emph{Transportation cost for Gaussian and other product measures.}
Geometric and Functional Analysis 6 (1996), 587--600.
\url{https://doi.org/10.1007/BF02249265}.
\end{thebibliography}
\end{document}
